\documentclass[11pt,a4paper]{article}

% Essential packages
\usepackage[utf8]{inputenc}
\usepackage[T1]{fontenc}
\usepackage{amsmath,amssymb,amsfonts}
\usepackage{physics}
\usepackage{braket}
\usepackage{geometry}
\usepackage{graphicx}
\usepackage{booktabs}
\usepackage{array}
\usepackage{hyperref}
\usepackage{xcolor}

\usepackage{multirow}

\usepackage{amsmath,amssymb,amsthm}
\usepackage{geometry}
\usepackage{booktabs}
\usepackage{array}
\usepackage{xcolor}
\usepackage{tcolorbox}
\usepackage{hyperref}
\usepackage{mathtools}
\usepackage{bm}
\usepackage{enumitem}




\newcommand{\Ms}{M_\odot}
\geometry{margin=2.5cm}

\definecolor{qgdblue}{RGB}{30,58,138}
\definecolor{qgdred}{RGB}{127,29,29}
\definecolor{qgdgreen}{RGB}{20,83,45}
\definecolor{boxbg}{RGB}{239,246,255}
\definecolor{boxborder}{RGB}{37,99,235}
\definecolor{greenbg}{RGB}{240,253,244}
\definecolor{redbg}{RGB}{254,242,242}

\newcommand{\Fp}{F_{\rm QGD}^{\rm dip}}
\newcommand{\Fq}{F_{\rm GR}^{\rm quad}}

\tcbuselibrary{theorems,skins,breakable}

\newtcbtheorem[number within=section]{theorem}{Theorem}{%
	enhanced,breakable,
	colback=boxbg,colframe=boxborder,
	fonttitle=\bfseries\color{qgdblue},
	theorem style=plain,separator sign none,
	attach boxed title to top left={yshift=-2mm,xshift=4mm},
	boxed title style={colback=boxborder}}{thm}

\newtcbtheorem[number within=section]{proposition}{Proposition}{%
	enhanced,breakable,
	colback=greenbg,colframe=qgdgreen!60,
	fonttitle=\bfseries\color{qgdgreen},
	theorem style=plain,separator sign none,
	attach boxed title to top left={yshift=-2mm,xshift=4mm},
	boxed title style={colback=qgdgreen!60}}{prop}

\newtcbtheorem[number within=section]{corollary}{Corollary}{%
	enhanced,breakable,
	colback=redbg,colframe=qgdred!50,
	fonttitle=\bfseries\color{qgdred},
	theorem style=plain,separator sign none,
	attach boxed title to top left={yshift=-2mm,xshift=4mm},
	boxed title style={colback=qgdred!50}}{cor}

\theoremstyle{definition}
\newtheorem{definition}{Definition}[section]
\theoremstyle{remark}
\newtheorem{remark}{Remark}[section]

% Custom commands
\newcommand{\boldr}{\mathbf{r}}
\newcommand{\boldv}{\mathbf{v}}
\newcommand{\bolda}{\mathbf{a}}
\newcommand{\bhat}[1]{\hat{\mathbf{#1}}}
\newcommand{\Sigt}{\Sigma_{\rm tot}}
\newcommand{\Mbox}{\Box_g}
\newcommand{\ETA}{\eta_{\mu\nu}}
\newcommand{\rs}[1]{r_{\rm s}^{(#1)}}
\newcommand{\crossterm}[2]{\sigma_{t}^{(#1)}\sigma_{t}^{(#2)}}
\newcommand{\sig}[2]{\sigma_{#1}^{(#2)}}
\newcommand{\lQ}{\ell_{Q}}
\newcommand{\boldx}{\mathbf{x}}
\newcommand{\Sigtot}{\Sigma_{\text{tot}}}
\newcommand{\bfa}{\mathbf{a}}

% Page layout
\geometry{
	a4paper,
	margin=1in,
	headheight=14pt
}

% Hyperref setup
\hypersetup{
	colorlinks=true,
	linkcolor=blue,
	citecolor=blue,
	urlcolor=blue,
	pdfauthor={},
	pdftitle={Quantum Gravity Dynamics: Emergent spacetime geometry from Dirac spinor fields}
}

% Custom commands

\title{Quantum Gravitational Dynamics:\\
	The Formulation of Quantum Gravity from the Dirac Equation}
\author{Romeo Matshaba}
\date{University of South Africa}

\begin{document}
	\maketitle
	
	\begin{abstract}
		We demonstrate that galaxy rotation curves across 175 SPARC galaxies (4,248 measurements) are reproduced with R²=0.908 and zero free parameters per galaxy by a quantum correction to the gravitational potential derived from the WKB phase of the Dirac wavefunction
	\end{abstract}
	
	\section{Introduction}
	
	The central question addressed in this paper is how classical gravity emerges from quantum mechanics.  The proposed answer rests on a specific mathematical structure: the wavefunction encodes gravitational dynamics, and current conservation provides the bridge to classical physics.
	
	The point of departure is the observation that gravity may be modelled as a spherical matter wave radiating to infinity.  Any such wave in three spatial dimensions carries a $1/r$ amplitude factor from geometry alone.  The wavefunction is therefore taken to have the form
	\begin{equation}
		\psi = \frac{\alpha_G}{r}\,u\,e^{iS/\hbar}
	\end{equation}
	giving $|\psi|^2 \propto 1/r^2$.  Combined with the current-conservation relation
	\begin{equation}
		|\psi|^2 \cdot p \cdot r^2 = C
	\end{equation}
	the $r^2$ factors cancel, yielding a momentum equation determined by geometry alone, without a plane-wave postulate.
	
	The analysis is carried out in \textit{flat Minkowski spacetime}; no curved background geometry is assumed.
	
	\bigskip
	\noindent\textbf{Note on the gravitational coupling constant.}
	Throughout this paper the single symbol $\alpha_G$ denotes the gravitational fine structure constant.  Its fundamental (complex) value, derived in Section~\ref{sec:newton} from phase-oscillation analysis, is
	\begin{equation}
		\boxed{\alpha_G = e^{i\pi/4}\sqrt{\frac{c\hbar}{2GMm}} = \frac{1+i}{\sqrt{2}}\sqrt{\frac{c\hbar}{2GMm}}}
		\label{eq:alphaG_def}
	\end{equation}
	so that
	\begin{equation}
		\alpha_G^2 = \frac{ic\hbar}{2GMm}, \qquad |\alpha_G|^2 = \frac{c\hbar}{2GMm}.
		\label{eq:alphaG_props}
	\end{equation}
	Because $\alpha_G$ is complex, it appears in the wavefunction and phase equations in its full form.  Wherever a \emph{real} quantity is required---probability densities, physical forces, normalization constants---one uses $|\alpha_G|^2$ or
	\begin{equation}
		\mathrm{Re}(\alpha_G) = \frac{1}{\sqrt{2}}\sqrt{\frac{c\hbar}{2GMm}} = \sqrt{\frac{c\hbar}{4GMm}}.
		\label{eq:alphaG_real}
	\end{equation}
	Dimensional check: $[\alpha_G^2] = [c\hbar/(GMm)] = 1$.~\checkmark
	
	%=============================================================
	\section{Microscopic Foundation}
	%=============================================================
	
	\subsection{Field Content and Dynamics}

	The theory is formulated in terms of a single microscopic degree of freedom: the Dirac spinor $\psi(x)\in\mathbb{C}^4$.  The microscopic action in flat Minkowski spacetime is
	\begin{equation}
	  \mathcal{S}_{\rm micro} = \int d^4x\!\left[
	    \bar{\psi}(i\gamma^\mu\partial_\mu - m)\psi
	    + G_{\rm NJL}(\bar{\psi}\psi)^2
	    - \tfrac{1}{4}F_{\mu\nu}F^{\mu\nu}
	    - e\bar{\psi}\gamma^\mu A_\mu\psi
	  \right]
	  \label{eq:S_micro_correct}
	\end{equation}
	where $\bar{\psi}=\psi^\dagger\gamma^0$, $\{\gamma^\mu,\gamma^\nu\}=2\eta^{\mu\nu}$ with $\eta=\mathrm{diag}(-1,+1,+1,+1)$, and $G_{\rm NJL}$ is the Nambu--Jona-Lasinio (NJL) four-fermion coupling constant.  The four terms are, respectively: the free Dirac kinetic term, a local four-fermion self-interaction, the Maxwell kinetic term, and the minimal electromagnetic coupling.

	The free-field subsector
	\begin{equation}
		S_D = \int d^4x\,\bar{\psi}(i\gamma^\mu\partial_\mu - m)\psi
		\label{eq:dirac_action}
	\end{equation}
	is the unique Lorentz-invariant, local, linear action for spin-$\tfrac{1}{2}$ fields with positive-definite energy, but it does not suffice as the microscopic starting point: coarse-graining $S_D$ alone generates only mass and wavefunction renormalisation, since no new interaction operators can arise from a free theory.  The pressure term $P(\bar\psi\psi)$, Maxwell kinetic term, and angular-momentum coupling $J^2\bar\psi\psi$ that appear in the gravitational effective Lagrangian all require the interacting terms present in $\mathcal{S}_{\rm micro}$.  Specifically, Wilsonian coarse-graining of $\mathcal{S}_{\rm micro}$ generates $P(\bar\psi\psi)$ from the four-Fermi term, $-\tfrac{1}{4}F_{\mu\nu}F^{\mu\nu}$ from the EM sector, the $J^2$ operator from its derivative expansion, and $\rho_\Lambda$ from zero-point fluctuations.  All coarse-graining results in the subsequent sections are derived from $\mathcal{S}_{\rm micro}$.

	\subsection{Stress-Energy Tensor}
	
	The symmetric stress-energy tensor:
	\begin{equation}
		T^{\mu\nu} = \frac{i}{4}\bar{\psi}\!\left(\gamma^\mu\overleftrightarrow{\partial}^\nu + \gamma^\nu\overleftrightarrow{\partial}^\mu\right)\!\psi
		\label{eq:stress_energy}
	\end{equation}
	where $\overleftrightarrow{\partial}=\overrightarrow{\partial}-\overleftarrow{\partial}$.  On-shell: $\partial_\mu T^{\mu\nu}=0$.
	
	This tensor provides the matter content that sources gravitational effects---even though we remain in flat spacetime.
	
	%=============================================================
	\section{Coarse-Graining and the Emergence of the Effective Theory}
	\label{sec:coarse}
	%=============================================================
	
	\subsection{Scale Hierarchy}
	
	The theory possesses three characteristic length scales.  Two of them are
	given at the outset:
	\begin{itemize}
		\item $\Lambda_{\mathrm{UV}}^{-1}$: the ultraviolet length scale at which
		the effective description breaks down.  Its physical identity is
		\emph{not assumed here}; it will be derived in
		Section~\ref{sec:UV_identification} as a consequence of the
		coarse-graining procedure and the structure of the loop integrals.
		
		\item $L_{\mathrm{grav}}$: the scale over which gravitational gradients
		are appreciable, i.e.\ the scale on which $\nabla T^{\mu\nu}/T^{\mu\nu}
		\sim 1/L_{\mathrm{grav}}$.
	\end{itemize}
	The coarse-graining scale $\ell$ is then any scale satisfying
	\begin{equation}
		\Lambda_{\mathrm{UV}}^{-1} \;\ll\; \ell \;\ll\; L_{\mathrm{grav}}.
		\label{eq:hierarchy}
	\end{equation}
	This hierarchy is a \emph{consistency requirement}, not an assumption: the
	Wilsonian procedure of Section~\ref{sec:mode_int} is self-consistent if and
	only if \eqref{eq:hierarchy} holds.  When $\ell \sim \Lambda_{\mathrm{UV}}^{-1}$,
	fast-mode fluctuations are not suppressed and the effective description fails;
	when $\ell \sim L_{\mathrm{grav}}$, the coarse-grained theory cannot resolve
	the gravitational gradients it is meant to describe.
	
		The identification $\Lambda_{\mathrm{UV}}^{-1} = \hbar/(mc)$ is a \emph{result},
		not an input.  It will emerge in Section~\ref{sec:UV_identification} from the
		requirement that the effective mass generated by the loop integrals matches the
		physical mass $m$.  We deliberately leave $\Lambda_{\mathrm{UV}}$ as a free
		UV scale throughout the present section to avoid any circularity.
	
	\subsection{Full Theory and Regularisation}

	The starting point for the Wilsonian coarse-graining is the path integral over $\mathcal{S}_{\rm micro}$~\eqref{eq:S_micro_correct} in flat Minkowski spacetime $(\mathcal{M}, \eta_{\mu\nu})$, signature $(-,+,+,+)$:
	\begin{equation}
		\mathcal{Z} \;=\;
		\int \mathcal{D}\bar{\psi}\,\mathcal{D}\psi\;
		\exp\!\left(\frac{i}{\hbar}
		\int d^4x\;
		\Bigl[\bar{\psi}\bigl(i\hbar\gamma^\mu\partial_\mu - mc\bigr)\psi
		+ G_{\rm NJL}(\bar{\psi}\psi)^2
		- \tfrac{1}{4}F_{\mu\nu}F^{\mu\nu}
		- e\bar{\psi}\gamma^\mu A_\mu\psi\Bigr]
		\right),
		\label{eq:path_integral}
	\end{equation}
	with a UV cutoff $\Lambda_{\mathrm{UV}}$ imposed by Pauli--Villars regularisation.  No assumption is made about the value or interpretation of $\Lambda_{\mathrm{UV}}$ at this stage.
	
	\subsection{Mode Decomposition and Integration over Fast Modes}
	\label{sec:mode_int}
	
	\subsubsection{Splitting the field}
	
	Introduce a sharp momentum cutoff at $\mu = 1/\ell$ and decompose:
	\begin{equation}
		\psi(\mathbf{x}) \;=\; \psi_{<}(\mathbf{x}) + \psi_{>}(\mathbf{x}),
		\label{eq:field_split}
	\end{equation}
	where
	\begin{align}
		\psi_{<}(\mathbf{x})
		&= \int_{|\mathbf{k}|<\mu}
		\frac{d^4k}{(2\pi)^4}\,\tilde{\psi}(\mathbf{k})\,e^{ik\cdot x},\\
		\psi_{>}(\mathbf{x})
		&= \int_{\mu\,<\,|\mathbf{k}|<\Lambda_{\mathrm{UV}}}
		\frac{d^4k}{(2\pi)^4}\,\tilde{\psi}(\mathbf{k})\,e^{ik\cdot x}.
	\end{align}
	The slow field $\psi_{<}$ encodes physics on scales $\geq \ell$; the fast field
	$\psi_{>}$ encodes physics between $\Lambda_{\mathrm{UV}}^{-1}$ and $\ell$.
	
	\subsubsection{Gaussian integration over fast modes}

	The fast-mode integral is evaluated by separating the action into a quadratic (free Dirac) part and the interaction terms.  The free part generates a Gaussian integral over $\psi_{>}$ that can be performed exactly; the NJL four-Fermi interaction and the EM coupling contribute via an expansion of $e^{iS_{\rm int}/\hbar}$ in powers of $G_{\rm NJL}$ and $e$ respectively, generating the loop diagrams that produce the effective operators listed below.  The result is
	\begin{equation}
		\mathcal{Z} \;=\;
		\int \mathcal{D}\bar{\psi}_{<}\,\mathcal{D}\psi_{<}\;
		\exp\!\left(\frac{i}{\hbar}S_{\mathrm{eff}}[\psi_{<}]\right),
		\label{eq:effective_pi}
	\end{equation}
	where the effective action is
	\begin{equation}
		S_{\mathrm{eff}}[\psi_{<}]
		\;=\; S_D[\psi_{<}]
		\;-\; i\hbar\,\mathrm{Tr}\ln
		\bigl(i\hbar\gamma^\mu\partial_\mu - mc - V[\psi_{<}]\bigr)
		\Big|_{>},
		\label{eq:Seff_exact}
	\end{equation}
	$|_{>}$ denotes restriction to the momentum shell $[\mu, \Lambda_{\mathrm{UV}}]$, and $V[\psi_{<}]$ encodes the interaction vertices of $\mathcal{S}_{\rm micro}$ with slow-mode external legs.  The $\mathrm{Tr}\ln$ generates all one-particle-irreducible (1PI) diagrams with fast-mode internal lines and slow-mode external legs.
	
	Generation of effective operators
		\label{prop:operators}
		Expanding $S_{\mathrm{eff}}$ in \eqref{eq:Seff_exact} in powers of $\psi_{<}$
		and retaining operators of mass dimension $\leq 4$ (marginal and relevant
		under the RG, cf.\ Section~\ref{sec:RG}), subject to Lorentz invariance,
		$\mathrm{U}(1)$ phase symmetry, and locality, yields
		\begin{multline}
			\mathcal{L}_{\mathrm{eff}}
			\;=\;
			\frac{i}{2}\bar{\psi}_{<}\gamma^\mu\overleftrightarrow{\partial}_\mu\psi_{<}
			\;-\; m_{\mathrm{eff}}\,\bar{\psi}_{<}\psi_{<}
			\;-\; P\!\left(\bar{\psi}_{<}\psi_{<}\right) \\
			-\; \frac{1}{4}F_{\mu\nu}F^{\mu\nu}
			\;-\; \frac{1}{2M^2}J_{\mu\nu}J^{\mu\nu}\bar{\psi}_{<}\psi_{<}
			\;-\; \rho_\Lambda,
			\label{eq:Leff}
		\end{multline}
		where $m_{\mathrm{eff}} = m + \delta m(\Lambda_{\mathrm{UV}},\mu)$ contains
		the mass renormalisation from fast-mode loops, with $\delta m$ depending
		explicitly on $\Lambda_{\mathrm{UV}}$.
	
		The loop expansion of \eqref{eq:Seff_exact} generates the following operators,
		each traceable to a specific diagram topology:
		\begin{enumerate}
			\item \textbf{Mass renormalisation.}
			The one-loop self-energy with two external slow legs shifts
			\begin{equation}
				\delta m \;=\; \frac{3e^2}{16\pi^2}\int_\mu^{\Lambda_{\mathrm{UV}}}
				\frac{d^4k}{(2\pi)^4}\,\frac{1}{k^2 - m^2c^2}
				\;\sim\; \frac{3e^2\Lambda_{\mathrm{UV}}^2}{16\pi^2 mc^2}.
				\label{eq:deltam}
			\end{equation}
			The physical mass is defined by the renormalisation condition
			$m_{\mathrm{eff}}|_{\mu\,=\,0} = m_{\mathrm{phys}}$, which fixes
			the counterterm and eliminates $\delta m$ from the final theory.
			The scale at which $\delta m$ becomes comparable to $m$ will identify
			$\Lambda_{\mathrm{UV}}$ (see Section~\ref{sec:UV_identification}).
			
			\item \textbf{Pressure.}
			The four-Fermi vertex $G_{\rm NJL}(\bar{\psi}\psi)^2$ present in $\mathcal{S}_{\rm micro}$ generates, via a four-external-leg slow-mode diagram with one fast internal loop, the term $-g_P(\bar{\psi}_{<}\psi_{<})^2$, which at low energy becomes the equation-of-state term $P(\bar{\psi}_{<}\psi_{<})$.
			
			\item \textbf{Electromagnetic coupling.}
			The minimal coupling $-e\bar{\psi}\gamma^\mu A_\mu\psi$ present in $\mathcal{S}_{\rm micro}$ is the unique $\mathrm{U}(1)$-covariant marginal deformation; integrating out fast charged modes generates the Maxwell kinetic term $-\tfrac{1}{4}F_{\mu\nu}F^{\mu\nu}$ in the effective theory.
			
			\item \textbf{Angular momentum / spin.}
			The dimension-6 operator
			$J_{\mu\nu}J^{\mu\nu}\bar{\psi}_{<}\psi_{<}/(2M^2)$ arises from
			box diagrams; it encodes frame-dragging effects in the stress-energy
			tensor.
			
			\item \textbf{Vacuum energy.}
			Zero-point fluctuations of fast modes contribute
			$\rho_\Lambda = \hbar\int_\mu^{\Lambda_{\mathrm{UV}}}
			d^4k/(2\pi)^4\cdot\omega_{\mathbf{k}}$, a constant shift whose
			IR value is set by matching to the observed cosmological constant.
			
			\item \textbf{No higher operators.}
			All operators of mass dimension $>4$ carry factors of
			$(\ell\,\Lambda_{\mathrm{UV}})^{-(d-4)} \ll 1$ under the hierarchy
			\eqref{eq:hierarchy} and are therefore suppressed. \hfill$\blacksquare$
		\end{enumerate}
	
	\subsection{Wilsonian Renormalisation Group Flow}
	\label{sec:RG}
	
	\subsubsection{RG transformation}
	
	Define the Wilsonian RG step as: (i) integrate out the momentum shell
	$|\mathbf{k}| \in [b^{-1}\mu, \mu]$, $b > 1$; (ii) rescale
	\begin{equation}
		x \;\to\; bx, \qquad
		\psi_{<} \;\to\; b^{-3/2}\psi_{<},
		\label{eq:rescaling}
	\end{equation}
	to restore the kinetic term to canonical normalisation.  Let $t = \ln b$.
	
	\subsubsection{Beta functions}
	
	The one-loop RG equations for the couplings in \eqref{eq:Leff} are:
	\begin{align}
		\frac{dm}{dt}
		&= m\left(1 + \gamma_m\right),
		\qquad \gamma_m = -\frac{3e^2}{16\pi^2}, \label{eq:beta_m}\\
		\frac{d g_P}{dt}
		&= g_P\!\left(1 - 2\gamma_\psi\right) + \frac{N_f g_P^2}{8\pi^2},
		\label{eq:beta_gP}\\
		\frac{d e^2}{dt}
		&= \frac{e^4}{6\pi^2}, \label{eq:beta_e}\\
		\frac{d(1/M^2)}{dt}
		&= -\frac{2}{M^2}\left(1 - \gamma_J\right),
		\label{eq:beta_M}\\
		\frac{d\rho_\Lambda}{dt}
		&= 4\rho_\Lambda + \frac{N_f\,\mu^4}{16\pi^2}.
		\label{eq:beta_Lambda}
	\end{align}
	Here $N_f$ is the number of fermion flavours and $\gamma_\psi$, $\gamma_J$
	are the anomalous dimensions of $\bar{\psi}\psi$ and $J_{\mu\nu}$
	respectively.
	
	IR fixed point and decoupling
		\label{prop:IR}
		In the limit $t \to \infty$ (i.e.\ $\ell \to L_{\mathrm{grav}}$):
		\begin{enumerate}
			\item The dimension-6 coupling $1/M^2 \to 0$: angular-momentum corrections
			are suppressed as $(\ell/L_{\mathrm{grav}})^2$.
			\item $\rho_\Lambda$ reaches its IR value fixed by cosmological matching.
			\item Quantum fluctuations average to zero:
			$\langle\psi_{>}\rangle = 0$, leaving classical effective fields.
			\item The physical mass $m_{\mathrm{phys}}$ is stable, with
			$\delta m$ absorbed by the renormalisation condition.
		\end{enumerate}
	
		The coupling $1/M^2$ has canonical dimension $-2$; under rescaling
		\eqref{eq:rescaling} it acquires a factor $b^{-2}$, so it is power-law
		irrelevant and flows to zero.  The remaining claims follow from the
		Appelquist--Carazzone decoupling theorem: fields with masses
		$\gg \mu$ decouple from low-energy physics, ensuring
		$\langle\psi_{>}\rangle = 0$ at scales $\ell \gg \Lambda_{\mathrm{UV}}^{-1}$.
		\hfill$\blacksquare$
	
	\subsection{The Effective Lagrangian}
	
	Effective Lagrangian from Wilsonian coarse-graining
		\label{thm:Leff}
		Let the interacting Dirac theory $\mathcal{S}_{\rm micro}$~\eqref{eq:S_micro_correct} be coarse-grained over a scale
		$\ell$ satisfying \eqref{eq:hierarchy}, with $\Lambda_{\mathrm{UV}}$ an
		unspecified UV scale.  Then the unique effective Lagrangian consistent with
		Lorentz invariance, $\mathrm{U}(1)$ symmetry, locality, and positive-definite
		energy, retaining only marginal and relevant operators, is
		\begin{equation}
			\boxed{
				\mathcal{L}_{\mathrm{eff}}
				\;=\;\frac{i}{2}\bar{\psi}\gamma^\mu\overleftrightarrow{\partial}_\mu\psi
				\;-\; m\bar{\psi}\psi
				\;-\; P(\bar{\psi}\psi)
				\;-\; \frac{1}{4}F_{\mu\nu}F^{\mu\nu}
				\;-\; \frac{1}{2M^2}J_{\mu\nu}J^{\mu\nu}\bar{\psi}\psi
				\;-\; \rho_\Lambda
			}
			\label{eq:Leff_final}
		\end{equation}
		(slow-mode subscripts dropped).  Each term has a precise gravitational
		interpretation:
		\begin{center}
			\begin{tabular}{ll}
				\hline
				\textbf{Term} & \textbf{Gravitational limit} \\
				\hline
				$m\bar{\psi}\psi$
				& Rest-mass density $\to$ Schwarzschild \\
				$P(\bar{\psi}\psi)$
				& Pressure $\to$ equation of state \\
				$\tfrac{1}{2M^2}J_{\mu\nu}J^{\mu\nu}\bar{\psi}\psi$
				& Angular momentum $\to$ Kerr frame-dragging \\
				$\tfrac{1}{4}F_{\mu\nu}F^{\mu\nu}$
				& EM energy $\to$ Reissner--Nordstr\"om \\
				$\rho_\Lambda$
				& Vacuum energy $\to$ cosmological constant \\
				\hline
			\end{tabular}
		\end{center}
		This correspondence is a \emph{result} of the coarse-graining, not a postulate.
	
		Proposition~\ref{prop:operators} establishes generation of all operators in
		\eqref{eq:Leff_final}.  Proposition~\ref{prop:IR} establishes that no further
		operators survive at IR scales.  Uniqueness follows from the exhaustive
		classification of Lorentz-invariant, $\mathrm{U}(1)$-symmetric, local
		bilinears and quadratics in $\psi$, $F_{\mu\nu}$, $J_{\mu\nu}$ of mass
		dimension $\leq 4$: every entry in \eqref{eq:Leff_final} appears in this
		list, and the list is finite. Higher-dimension operators are suppressed by
		$(\Lambda_{\mathrm{UV}}^{-1}/\ell)^{d-4} \ll 1$ under \eqref{eq:hierarchy}.
		\hfill$\blacksquare$
	
	\subsection{Connection to the Phase Gradient and Newton's Law}
	\label{sec:coarse:F}
	
	After coarse-graining, apply the WKB ansatz to the slow mode:
	\begin{equation}
		\psi_{<}(\mathbf{x})
		\;=\; R(\mathbf{x})\,e^{iS(\mathbf{x})/\hbar}
	\end{equation}
	where $R$ varies on scales $\sim L_{\mathrm{grav}}$ and $S$ oscillates on
	scales $\gg \ell$.  The phase gradient $\sigma_\mu$ is the effective classical
	variable from which the metric is reconstructed.
	
	Applying Noether's theorem to $\mathcal{L}_{\mathrm{eff}}$ and averaging over
	$\ell$ yields the coarse-grained stress-energy tensor:
	\begin{equation}
		\boxed{
			\langle T^{\mu\nu}\rangle
			\;=\; \rho c^2 u^\mu u^\nu
			\;+\; Pg^{\mu\nu}
			\;+\; \langle\tau^{\mu\nu}_{\mathrm{spin}}\rangle
			\;+\; \langle\tau^{\mu\nu}_{\mathrm{EM}}\rangle
			\;+\; \rho_\Lambda g^{\mu\nu}.
		}
		\label{eq:Tmunu_full}
	\end{equation}
	The result~\eqref{eq:Tmunu_full} has the structure of the Einstein stress-energy tensor; each term arises from the corresponding term in the effective Lagrangian via coarse-graining rather than being postulated independently.
	
	The gravitational force function then takes the form $F(r) = \Omega/P(r)$
	where $P(r) = i\sinh(2\alpha r)$ emerges from the phase-gradient equation;
	Newton's law is the $\alpha r \to 0$ limit (see Section~\ref{sec:newton} and
	Section~\ref{sec:UV_identification} where $\alpha$ --- and hence
	$\Lambda_{\mathrm{UV}}^{-1}$ --- is identified).
	
	\subsection{Effective Lagrangian}
	\label{sec:eff_lagrangian_summary}

	% NOTE (C8): This subsection is retained as a navigational anchor but the
	% Lagrangian is not repeated here to avoid duplication with Theorem~\ref{thm:Leff}
	% above, which states and derives the identical expression eq.~\eqref{eq:Leff_final}
	% with its full physical interpretation table.  See that theorem for the complete
	% statement.  The label \texttt{eq:eff\_lagrangian} below aliases the same content.
	\newcommand{\efflagalias}{\eqref{eq:Leff_final}}%

	The effective Lagrangian derived in Theorem~\ref{thm:Leff}, Eq.~\eqref{eq:Leff_final},
	is the unique result consistent with Lorentz invariance, U(1) symmetry, locality,
	and positive-definite energy.  To summarise its physical content: each term maps to
	a specific classical gravitational effect in the macroscopic limit---rest-mass density
	to Schwarzschild, pressure to equation of state, angular momentum to Kerr
	frame-dragging, electromagnetic energy to Reissner--Nordstr\"om, and vacuum energy to
	the cosmological constant.  The derivation and full table are given in
	Theorem~\ref{thm:Leff}.
	
	%=============================================================
	\section{The Complete Action}
	%=============================================================
	
	\begin{equation}
		\boxed{
			S = \int d^4x \left[\frac{i}{2}\bar{\psi}\gamma^\mu\overleftrightarrow{\partial}_\mu\psi - m\bar{\psi}\psi - P(\bar{\psi}\psi) - \frac{1}{2M^2}J_{\mu\nu}J^{\mu\nu}\bar{\psi}\psi - \frac{1}{4}F_{\mu\nu}F^{\mu\nu} - \rho_\Lambda\right]
		}
		\label{eq:complete_action}
	\end{equation}
	
	where:
	\begin{itemize}
		\item $\bar{\psi}=\psi^\dagger\gamma^0$,\quad $\{\gamma^\mu,\gamma^\nu\}=2\eta^{\mu\nu}$
		\item $\overleftrightarrow{\partial}_\mu = \overrightarrow{\partial}_\mu - \overleftarrow{\partial}_\mu$
		\item $J_{\mu\nu} = x_\mu p_\nu - x_\nu p_\mu$ (angular momentum tensor)
		\item $F_{\mu\nu} = \partial_\mu A_\nu - \partial_\nu A_\mu$ (electromagnetic field tensor)
	\end{itemize}
	
	%=============================================================
	\section{Stress-Energy Tensor from Action}
	%=============================================================
	
	\subsection{Noether's Theorem}
	
	From spacetime translation invariance:
	\begin{equation}
		T^{\mu\nu} = \frac{\partial\mathcal{L}}{\partial(\partial_\mu\psi)}\partial^\nu\psi + \frac{\partial\mathcal{L}}{\partial(\partial_\mu\bar{\psi})}\partial^\nu\bar{\psi} - g^{\mu\nu}\mathcal{L}
	\end{equation}
	
	After symmetrization and coarse-graining over scale $\ell$:
	\begin{equation}
		\boxed{\langle T^{\mu\nu}\rangle = \rho(x)c^2 u^\mu u^\nu + P(x)g^{\mu\nu} + \langle\tau^{\mu\nu}_{\mathrm{spin}}\rangle + \langle\tau^{\mu\nu}_{\mathrm{EM}}\rangle + \rho_\Lambda g^{\mu\nu}}
		\label{eq:stress_energy_full}
	\end{equation}
	
	The result has the structure of the Einstein stress-energy tensor.
	
	\subsection{Component Origins}
	
	Each stress-energy term traces to an action term:
	\begin{align}
		\rho c^2 u^\mu u^\nu &\leftarrow \tfrac{i}{2}\bar{\psi}\gamma^\mu\overleftrightarrow{\partial}_\mu\psi - m\bar{\psi}\psi \\
		Pg^{\mu\nu} &\leftarrow -P(\bar{\psi}\psi) \\
		\tau^{\mu\nu}_{\mathrm{spin}} &\leftarrow -\tfrac{1}{2M^2}J_{\mu\nu}J^{\mu\nu}\bar{\psi}\psi \\
		\tau^{\mu\nu}_{\mathrm{EM}} &\leftarrow -\tfrac{1}{4}F_{\mu\nu}F^{\mu\nu} \\
		\rho_\Lambda g^{\mu\nu} &\leftarrow -\rho_\Lambda
	\end{align}
	
	%=============================================================
	\section{Modified Dirac Equation and Energy Denominators}
	%=============================================================
	
	\subsection{Equation of Motion}
	
	Varying the action with respect to $\bar{\psi}$:
	\begin{equation}
		i\gamma^\mu\partial_\mu\psi = \left[m + \frac{\partial P}{\partial(\bar{\psi}\psi)} + \frac{1}{2M^2}J_{\mu\nu}J^{\mu\nu}\right]\psi
	\end{equation}
	
	For stationary states $\psi = u(p)e^{-iS/\hbar}$:
	\begin{equation}
		[\gamma^\mu p_\mu - mc\gamma^0]\psi = V_{\mathrm{eff}}\gamma^0\psi
	\end{equation}
	
	\subsection{The Effective Hamiltonian}
	
	All stress-energy contributions enter through:
	\begin{equation}
		\boxed{H = \int\rho c^2\,dV - \int P\,dV - \frac{J^2}{2mr^2} + V_{EM} + \rho_\Lambda}
		\label{eq:hamiltonian}
	\end{equation}
	
	where each term originates from the effective Lagrangian:
	\begin{itemize}
		\item $\int\rho c^2\,dV$: from $\frac{i}{2}\bar{\psi}\gamma^\mu\overleftrightarrow{\partial}_\mu\psi$ (source mass-energy)
		\item $-\int P\,dV$: from $-P(\bar{\psi}\psi)$ (pressure work)
		\item $-J^2/(2mr^2)$: from $-\frac{1}{2M^2}J_{\mu\nu}J^{\mu\nu}\bar{\psi}\psi$ (frame-dragging)
		\item $V_{EM}$: from $-\frac{1}{4}F_{\mu\nu}F^{\mu\nu}$ (electromagnetic potential)
		\item $\rho_\Lambda$: from $-\rho_\Lambda$ (vacuum energy)
	\end{itemize}
	
	\subsection{Energy Denominators}
	
	This defines the energy denominators:
	\begin{align}
		\Delta  &= E + mc^2 + H \quad\text{(particle states)}    \label{eq:Delta}\\
		\Delta' &= E - mc^2 - H \quad\text{(antiparticle states)} \label{eq:Delta_prime}
	\end{align}
	
	\textbf{Explicit form with all contributions:}
	\begin{equation}
		\boxed{
			\Delta = E + \underbrace{mc^2}_{\text{test}} + \underbrace{\int\rho c^2\,dV}_{\text{source}} - \underbrace{\int P\,dV}_{\text{pressure}} - \underbrace{\frac{J^2}{2mr^2}}_{\text{spin}} + \underbrace{V_{EM}}_{\text{EM}} + \underbrace{\rho_\Lambda}_{\text{vacuum}}
		}
		\label{eq:Delta_full}
	\end{equation}
	
	%=============================================================
	\section{The Quantum-Classical Bridge}
	%=============================================================
	
	\subsection{Macroscopic Coherent States and the Spherical Wave Ansatz}
	
	Gravity is treated as a \textbf{spherical matter wave radiating to infinity}.  Any wave radiating from a point source in three spatial dimensions must spread over a growing spherical surface of area $4\pi r^2$.  Flux conservation therefore demands that the wavefunction amplitude falls as $1/r$.  A macroscopically coherent gravitational state must take the form:
	\begin{equation}
		\psi = \frac{\alpha_G}{r}\,e^{iS(x)/\hbar}\,u
		\label{eq:coherent_state}
	\end{equation}
	where:
	\begin{itemize}
		\item $\alpha_G/r$: amplitude --- the $1/r$ is \textbf{geometric}, mandatory for any spherical wave in 3D
		\item $S(x)$: rapidly-varying phase (the action)
		\item $u$: slowly-varying four-component spinor envelope
	\end{itemize}
	
	\textbf{Dimensional note.}  Previously the wavefunction was written $\psi = \alpha_G \cdot u \cdot e^{iS/\hbar}$ with $\alpha_G$ carrying no $r$-dependence.  This was the source of the earlier inconsistency: $\alpha_G$ was dimensionless but $\psi$ requires dimensions $L^{-3/2}$, so the spatial structure was implicitly hidden in $u$.  Explicitly writing $1/r$ restores dimensional consistency and follows from the spherical-wave geometry.
	
	\subsection{Probability Current}
	
	The Dirac current with the spherical wave ansatz:
	\begin{equation}
		j^\mu = \bar{\psi}\gamma^\mu\psi = \frac{|\alpha_G|^2}{r^2}\,\bar{u}\gamma^\mu u
		\label{eq:current}
	\end{equation}
	
	\textit{Note:} The phase $e^{iS/\hbar}$ cancels in $\bar{\psi}\psi = \psi^\dagger\gamma^0\psi$.
	
	\textbf{Non-relativistic limit:} For spinor aligned with momentum:
	\begin{equation}
		j^0 \approx \frac{|\alpha_G|^2}{r^2}, \qquad \mathbf{j} \approx \frac{|\alpha_G|^2}{r^2}\frac{\nabla S}{m}
		\label{eq:current_NR}
	\end{equation}
	
	The total outward probability flux through any sphere of radius $r$ is:
	\begin{equation}
		\Phi = 4\pi r^2 \cdot j^r = 4\pi|\alpha_G|^2 \cdot p = \text{const}
	\end{equation}
	confirming that $|\alpha_G|^2 \sim \text{kg/s}$ is a conserved flux, as required.
	
	\subsection{Current Conservation: The Central Equation}
	
	From $\partial_\mu j^\mu = 0$, for stationary states:
	\begin{equation}
		\boxed{\nabla\cdot\!\left(|\psi|^2\nabla S\right) = 0}
		\label{eq:continuity}
	\end{equation}
	
	\subsection{Spherical Symmetry: The Exact Bridge}
	
	For spherically symmetric configurations with $\nabla S = p(r)\hat{r}$:
	\begin{equation}
		\frac{1}{r^2}\frac{\partial}{\partial r}\!\left(r^2|\psi|^2 p(r)\right) = 0
	\end{equation}
	
	Integrating gives the \textbf{exact} conservation law:
	\begin{equation}
		\boxed{|\psi|^2 \cdot p(r) \cdot r^2 = C}
		\label{eq:exact_bridge}
	\end{equation}
	
	Substituting $|\psi|^2 = (|\alpha_G|^2/r^2)(1 + p^2c^2/\Delta^2)$ from the spherical wave ansatz:
	\begin{equation}
		\frac{|\alpha_G|^2}{r^2}\!\left(1 + \frac{p^2}{\Delta^2}\right)\cdot p \cdot r^2 = C
	\end{equation}
	
	\textbf{The $r^2$ cancels exactly from geometry:}
	\begin{equation}
		\boxed{|\alpha_G|^2\!\left(p + \frac{p^3}{\Delta^2}\right) = C}
		\label{eq:cubic_from_geometry}
	\end{equation}
	
	Equation~\eqref{eq:cubic_from_geometry} is the quantum-classical bridge.  The $r^2$ cancellation follows from the spherical-wave ansatz and is exact; no plane-wave approximation is introduced.  The same geometric factor that produces Newton's $1/r^2$ law is responsible for the cancellation.
	
	The constant $C$ will be determined through the complete system of equations.

	\begin{remark}[The bridge holds for all $r$; the closing equation changes with regime]
	\label{rem:regimes}
	The current-conservation bridge $|\psi|^2\cdot p\cdot r^2 = C$ is \emph{exact} at every
	length scale---it follows from $\partial_\mu j^\mu = 0$ alone and requires no
	approximation.  It provides one equation in two unknowns $\{p(r),\,|\psi|(r)\}$
	and must be paired with a second, dynamical equation.  That second equation changes
	with the physical regime:
	\begin{center}
	\renewcommand{\arraystretch}{1.4}
	\begin{tabular}{@{}p{3cm}p{4.2cm}p{6cm}@{}}
	\hline
	\textbf{Regime} & \textbf{Dynamical equation} & \textbf{Result} \\
	\hline
	$r\sim\lambda_C$ (quantum gravity) &
	  Phase expansion of $e^{2ipr/\hbar}$ + gravitational matching &
	  Newton's law \emph{and} quantum corrections $F\propto(\lambda_C/r)^{2n}$ \\[4pt]
	$r\gg\lambda_C$, classical &
	  QGD field equation $\square_g\sigma=\text{source}$, \newline solution $\sigma_t=\sqrt{r_s/r}$ &
	  Newton's law (0PN), exact for all macroscopic $r$ \\[4pt]
	$r\gg\lambda_C$, relativistic &
	  QGD propagator loop integrals \newline $\Delta_{\rm QGD}(k)=1/k^2-1/(k^2+m_Q^2)$ &
	  Full PN potential $A(u;\nu)$ to all orders \\
	\hline
	\end{tabular}
	\end{center}
	The bridge $|\psi|^2\cdot p\cdot r^2=C$ is therefore the \textbf{universal skeleton}
	spanning all three regimes and all distance scales.  Demonstrating this unified
	coverage explicitly is a goal of the ongoing QGD programme.
	\end{remark}
	
	%=============================================================
	\section{Complete Four-Component Wavefunction Solution}
	%=============================================================
	
	\subsection{Eigenvalue Problem with Full Stress-Energy}
	
	In Pauli-Dirac representation $u=(u_A,u_B)^T$:
	\begin{equation}
		\begin{pmatrix} E-mc^2-H & -c\boldsymbol{\sigma}\cdot\mathbf{p} \\ c\boldsymbol{\sigma}\cdot\mathbf{p} & -(E+mc^2+H) \end{pmatrix}
		\begin{pmatrix} u_A \\ u_B \end{pmatrix} = 0
		\label{eq:coupled_system}
	\end{equation}
	where $H$ contains all stress-energy contributions as defined in Eq.~\eqref{eq:hamiltonian}.
	
	\subsection{General Solution with All Stress-Energy Terms}
	
	The most general solution, with $\alpha_G$ as the universal prefactor, is:
	
	\begin{equation}
		\psi = \frac{\alpha_G}{r}\Bigg[\psi_0\begin{bmatrix}1\\0\\\sigma_z\\\sigma_x+i\sigma_y\end{bmatrix}e^{-iS/\hbar}
		+ \psi_1\begin{bmatrix}0\\1\\\sigma_x-i\sigma_y\\-\sigma_z\end{bmatrix}e^{-iS/\hbar}
		+ \psi_2\begin{bmatrix}\sigma_z'\\\sigma_x'+i\sigma_y'\\1\\0\end{bmatrix}e^{+iS/\hbar}
		+ \psi_3\begin{bmatrix}\sigma_x'-i\sigma_y'\\-\sigma_z'\\0\\1\end{bmatrix}e^{+iS/\hbar}\Bigg]
		\label{eq:complete_solution}
	\end{equation}
	
	where the \textbf{graviton scalars} with minimal electromagnetic coupling are:
	\begin{align}
		\sigma_i  &= \frac{(p_i - eA_i/c)c}{\Delta}  \quad\text{(particle states)}    \label{eq:sigma}\\
		\sigma_i' &= \frac{(p_i - eA_i/c)c}{\Delta'} \quad\text{(antiparticle states)} \label{eq:sigma_prime}
	\end{align}
	and $\Delta$, $\Delta'$ contain all stress-energy contributions from Eq.~\eqref{eq:Delta_full}.
	
	\textbf{Note on reality:}  $\alpha_G$ is complex; in all probability-density computations we use $|\alpha_G|^2 = c\hbar/(2GMm)$, not $\alpha_G^2$.
	
	\subsection{Unified Form (Single-Component, Particle-Only)}
	
	\begin{equation}
		\boxed{
			\psi = \frac{\alpha_G}{r}\begin{bmatrix}
				1 \\ 0 \\
				\dfrac{(p_z-eA_z/c)c}{\Delta_{\mathrm{full}}} \\[6pt]
				\dfrac{[(p_x-eA_x/c)+i(p_y-eA_y/c)]c}{\Delta_{\mathrm{full}}}
			\end{bmatrix}e^{-iS/\hbar}
		}
	\end{equation}
	
	where:
	\begin{equation}
		\boxed{
			\Delta_{\mathrm{full}} = E + mc^2 + \int\rho c^2\,dV - \int P\,dV - \frac{J^2}{2mr^2} + V_{EM} + \rho_\Lambda
		}
		\label{eq:Delta_unified}
	\end{equation}
	
	All five stress-energy contributions ($\rho,P,J,F,\Lambda$) enter through the same denominator $\Delta_{\mathrm{full}}$.
	
	%=============================================================
	\section{Special Cases: Classical Limits}
	%=============================================================
	
	\subsection{Pure Schwarzschild ($P=J=F=\Lambda=0$)}
	
	\begin{equation}
		\Delta_{\mathrm{Schw}} = E + mc^2 + \int\rho c^2\,dV
	\end{equation}
	
	\begin{equation}
		\psi_{\mathrm{Schw}} = \frac{\alpha_G}{r}\begin{bmatrix}1\\0\\ p_zc/\Delta_{\mathrm{Schw}}\\ (p_x+ip_y)c/\Delta_{\mathrm{Schw}}\end{bmatrix}e^{-iS/\hbar}
	\end{equation}
	
	Stress-energy: $T^{\mu\nu} = \rho c^2 u^\mu u^\nu$
	
	\subsection{Kerr: Rotating Black Hole ($P=F=\Lambda=0$, $J\neq0$)}
	
	\begin{equation}
		\Delta_{\mathrm{Kerr}} = E + mc^2 + \int\rho c^2\,dV - \frac{J^2}{2mr^2}
	\end{equation}
	
	\begin{equation}
		\psi_{\mathrm{Kerr}} = \frac{\alpha_G}{r}\begin{bmatrix}1\\0\\ p_zc/\Delta_{\mathrm{Kerr}}\\ (p_x+ip_y)c/\Delta_{\mathrm{Kerr}}\end{bmatrix}e^{-iS/\hbar}
	\end{equation}
	
	Stress-energy: $T^{\mu\nu} = \rho c^2 u^\mu u^\nu + \tau^{\mu\nu}_{\mathrm{spin}}$
	
	\subsection{Reissner-Nordstr\"om: Charged Black Hole ($P=J=\Lambda=0$, $F\neq0$)}
	
	\begin{equation}
		\Delta_{\mathrm{RN}} = E + mc^2 + \int\rho c^2\,dV + V_{EM}
	\end{equation}
	
	\begin{equation}
		\psi_{\mathrm{RN}} = \frac{\alpha_G}{r}\begin{bmatrix}1\\0\\ (p_z-eA_z/c)c/\Delta_{\mathrm{RN}}\\ [(p_x-eA_x/c)+i(p_y-eA_y/c)]c/\Delta_{\mathrm{RN}}\end{bmatrix}e^{-iS/\hbar}
	\end{equation}
	
	Stress-energy: $T^{\mu\nu} = \rho c^2 u^\mu u^\nu + \tau^{\mu\nu}_{\mathrm{EM}}$
	
	\subsection{Perfect Fluid with Cosmological Constant ($J=F=0$)}
	
	\begin{equation}
		\Delta_{\mathrm{fluid}} = E + mc^2 + \int\rho c^2\,dV - \int P\,dV + \rho_\Lambda
	\end{equation}
	
	\begin{equation}
		\psi_{\mathrm{fluid}} = \frac{\alpha_G}{r}\begin{bmatrix}1\\0\\ p_zc/\Delta_{\mathrm{fluid}}\\ (p_x+ip_y)c/\Delta_{\mathrm{fluid}}\end{bmatrix}e^{-iS/\hbar}
	\end{equation}
	
	Stress-energy: $T^{\mu\nu} = \rho c^2 u^\mu u^\nu + Pg^{\mu\nu} + \rho_\Lambda g^{\mu\nu}$
	
	\subsection{Most General Case}
	
	\begin{equation}
		\boxed{
			\psi_{\mathrm{general}} = \frac{\alpha_G}{r}\begin{bmatrix}
				1 \\ 0 \\
				\dfrac{(p_z-eA_z/c)c}{E+mc^2+\int\rho c^2\,dV-\int P\,dV-\frac{J^2}{2mr^2}+V_{EM}+\rho_\Lambda} \\[10pt]
				\dfrac{[(p_x-eA_x/c)+i(p_y-eA_y/c)]c}{E+mc^2+\int\rho c^2\,dV-\int P\,dV-\frac{J^2}{2mr^2}+V_{EM}+\rho_\Lambda}
			\end{bmatrix}e^{-iS/\hbar}
		}
	\end{equation}
	
	%=============================================================
	\section{Connecting Quantum and Classical}
	%=============================================================
	
	\subsection{Probability Density}
	
	With the spherical wave ansatz $\psi = (\alpha_G/r)\,u\,e^{-iS/\hbar}$:
	\begin{align}
		|\psi|^2 &= \psi^\dagger\psi
		= \frac{|\alpha_G|^2}{r^2}\begin{bmatrix}1&0&\frac{p_zc}{\Delta}&\frac{(p_x+ip_y)c}{\Delta}\end{bmatrix}^\dagger
		\begin{bmatrix}1\\0\\\frac{p_zc}{\Delta}\\\frac{(p_x+ip_y)c}{\Delta}\end{bmatrix} \\
		&= \frac{|\alpha_G|^2}{r^2}\!\left(1 + \frac{|p|^2 c^2}{\Delta^2}\right)
	\end{align}
	
	Therefore:
	\begin{equation}
		\boxed{|\psi|^2 = \frac{|\alpha_G|^2}{r^2}\!\left(1 + \frac{|p|^2 c^2}{\Delta^2}\right)}
		\label{eq:prob_density}
	\end{equation}
	
	where $\Delta$ contains all stress-energy contributions.
	
	\subsection{Applying the Exact Bridge}
	
	The exact conservation law (Eq.~\eqref{eq:exact_bridge}) states $|\psi|^2 \cdot p \cdot r^2 = C$.
	
	Substituting Eq.~\eqref{eq:prob_density}:
	\begin{equation}
		\frac{|\alpha_G|^2}{r^2}\!\left(1+\frac{|p|^2 c^2}{\Delta^2}\right)\cdot|p|\cdot r^2 = C
		\label{eq:momentum_equation}
	\end{equation}
	
	\textbf{The $r^2$ cancels exactly --- a geometric consequence of the spherical wave:}
	\begin{equation}
		|\alpha_G|^2|p| + |\alpha_G|^2\frac{c^2|p|^3}{\Delta^2} = C
	\end{equation}
	
	\begin{equation}
		\boxed{|p| + \frac{c^2|p|^3}{\Delta^2} = \frac{C}{|\alpha_G|^2}}
		\label{eq:cubic}
	\end{equation}
	
	\textbf{This cubic equation determines the momentum structure from the wavefunction.}  The cancellation is exact, not an approximation.
	
	\subsection{Non-Relativistic Limit}
	
	For $|p|\ll mc$, $\Delta\approx mc^2$, the cubic term is negligible:
	\begin{equation}
		\boxed{|p| \approx \frac{C}{|\alpha_G|^2}}
		\label{eq:nr_momentum}
	\end{equation}
	
	This connects the normalization constants to momentum.
	
	%=============================================================
	\section{Derivation of Newton's Law and Determination of $\alpha_G$}
	\label{sec:newton}
	%=============================================================
	
	\subsection{Force from Wavefunction Phase Structure}
	
	The momentum and energy fields in the weak-field limit are:
	\begin{equation}
		p(r) = \frac{p_0}{\alpha_G^2}\,e^{-2ipr/\hbar}, \qquad
		E(r) = \frac{p_0 c}{\alpha_G^2\,e^{2ipr/\hbar}}.
	\end{equation}
	
	\subsection{Phase Expansion}
	
	Taylor-expanding with $p_0=mc$:
	\begin{equation}
		e^{2imcr/\hbar} = 1 + \frac{2imcr}{\hbar} - \frac{2m^2c^2r^2}{\hbar^2} - i\frac{4m^3c^3r^3}{3\hbar^3} + \cdots
	\end{equation}
	
	\subsection{Leading-Term Analysis}
	
	Using the first non-trivial term:
	\begin{equation}
		E(r) \approx \frac{mc^2}{\alpha_G^2\cdot 2imcr/\hbar} = \frac{c\hbar}{2i\,\alpha_G^2\,r}.
	\end{equation}
	
	\begin{tcolorbox}[enhanced,breakable,
	  colback=white,colframe=black,
	  title={\textbf{Key step: this is where the free-Dirac theory couples to gravity}},
	  fonttitle=\bfseries]
	Up to this point $\alpha_G$ is an \emph{undetermined} normalization constant of the
	free wavefunction.  It becomes the gravitational coupling constant only when we
	\textbf{match the phase-field energy to the Newtonian gravitational potential} of a
	source of mass $M$:
	\begin{equation}
	  \underbrace{\frac{c\hbar}{2i\,\alpha_G^2\,r}}_{\text{phase-field energy (free QM)}}
	  \;=\;
	  \underbrace{-\frac{GMm}{r}}_{\text{gravitational potential (input)}}.
	  \label{eq:coupling_step}
	\end{equation}
	This matching is the \emph{gravitational coupling step}: without it, $\alpha_G$
	remains a free parameter and no gravitational physics is predicted.
	After this step, $G$, $M$, and $m$ enter $\alpha_G$ and all subsequent quantities.

	\medskip
	\noindent\textbf{Note on the logical status of this step.}
	The Newtonian potential $V = -GMm/r$ on the right-hand side is an \emph{external
	physical input}, not derived from within QGD at this stage.  It introduces Newton's
	constant $G$, the source mass $M$, and the test mass $m$ as free parameters fixed
	by experiment.  What QGD \emph{does} derive from this matching is:
	\begin{enumerate}
	  \item The specific form $\alpha_G^2 = ic\hbar/(2GMm)$, which fixes the wavefunction
	    normalization in terms of the gravitational masses.
	  \item The $1/r^2$ falloff of the force (Section~\ref{sec:newton}.\ref{subsec:newton_force}),
	    which follows from $F = -dE/dr$ applied to $E \propto 1/r$ --- a consequence of
	    the spherical-wave geometry, not an additional input.
	  \item All quantum corrections beyond Newton (Section~\ref{sec:UV_identification}),
	    which are genuine predictions with no classical analogue.
	\end{enumerate}
	The recovery of $F = GMm/r^2$ in Section~\ref{sec:newton}.\ref{subsec:newton_force}
	therefore demonstrates \emph{self-consistency} of QGD with Newton's law, not an
	independent derivation of it.  The fundamental assumption is that classical
	Newtonian gravity holds at macroscopic scales; QGD provides the quantum-mechanical
	framework in which this is realised and extended.
	\end{tcolorbox}

	For a static gravitational source, $E_{\mathrm{field}} = V = -GMm/r$:
	\begin{equation}
		\frac{c\hbar}{2i\,\alpha_G^2\,r} = -\frac{GMm}{r}.
	\end{equation}
	
	Solving:
	\begin{equation}
		\boxed{\alpha_G^2 = \frac{ic\hbar}{2GMm}}
		\label{eq:alphaG2}
	\end{equation}
	
	\subsection{The Complex Gravitational Coupling}
	
	Taking the square root using $\sqrt{i}=e^{i\pi/4}=(1+i)/\sqrt{2}$:
	\begin{equation}
		\boxed{\alpha_G = e^{i\pi/4}\sqrt{\frac{c\hbar}{2GMm}} = \frac{1+i}{\sqrt{2}}\sqrt{\frac{c\hbar}{2GMm}}}
		\label{eq:alphaG_full}
	\end{equation}
	
	Verification:
	\begin{equation}
		\alpha_G^2 = \frac{(1+i)^2}{2}\cdot\frac{c\hbar}{2GMm} = \frac{2i}{2}\cdot\frac{c\hbar}{2GMm} = \frac{ic\hbar}{2GMm} \quad\checkmark
	\end{equation}
	
	Key derived quantities:
	\begin{align}
		|\alpha_G|^2       &= \frac{c\hbar}{2GMm}                             \label{eq:mod_alphaG} \\
		\mathrm{Re}(\alpha_G) &= \sqrt{\frac{c\hbar}{4GMm}}                  \label{eq:re_alphaG}
	\end{align}
	
	\subsection{Newton's Force from the Leading Term}
	\label{subsec:newton_force}
	
	\begin{equation}
		F_1 = -\frac{dE}{dr} = \frac{c\hbar}{2i\,\alpha_G^2\,r^2}.
	\end{equation}
	
	Substituting $\alpha_G^2 = ic\hbar/(2GMm)$:
	\begin{equation}
		F_1 = \frac{c\hbar}{2ir^2}\cdot\frac{2GMm}{ic\hbar} = \frac{GMm}{r^2} \quad\checkmark
	\end{equation}
	
	Newton's inverse-square law is recovered exactly.
	
	%=============================================================
	\section{Determination of the Normalization Constant $C$}
	%=============================================================
	
	\subsection{Classical Momentum at the Gravitational Bohr Radius}
	
	For a particle with zero total energy in a Newtonian potential:
	\begin{equation}
		p(r) = \sqrt{\frac{2GMm^2}{r}}.
	\end{equation}
	
	The gravitational Bohr radius $a_0 = \hbar^2/(GMm^2)$ gives:
	\begin{equation}
		p(a_0) = \sqrt{\frac{2GMm^2\cdot GMm^2}{\hbar^2}} = \frac{\sqrt{2}\,GMm^2}{\hbar}.
	\end{equation}
	
	\subsection{Value of $C$}
	
	From Eq.~\eqref{eq:nr_momentum}, $C = |\alpha_G|^2\cdot p(a_0)$:
	\begin{equation}
		C = \frac{c\hbar}{2GMm}\cdot\frac{\sqrt{2}\,GMm^2}{\hbar} = \frac{\sqrt{2}\,cm\cdot m}{2m} = \frac{mc}{\sqrt{2}}.
	\end{equation}
	
	\begin{equation}
		\boxed{C = \frac{mc}{\sqrt{2}}}
		\label{eq:C_final}
	\end{equation}
	
	\textbf{Remarks:}  $C$ has dimensions of momentum~\checkmark.  The flux scale is set purely by the test particle's rest momentum; source-mass dependence enters through $\alpha_G$ and $\Delta$.
	
	\subsection{The Cubic Momentum Equation}
	
	Substituting $C$ and $|\alpha_G|^2$ into Eq.~\eqref{eq:cubic}:
	\begin{equation}
		\frac{C}{|\alpha_G|^2} = \frac{mc/\sqrt{2}}{c\hbar/(2GMm)} = \frac{mc}{\sqrt{2}}\cdot\frac{2GMm}{c\hbar} = \frac{\sqrt{2}\,GMm^2}{\hbar} \equiv \mathcal{P}.
	\end{equation}
	
	\begin{equation}
		\boxed{|p| + \frac{c^2|p|^3}{\Delta^2} = \mathcal{P} = \frac{\sqrt{2}\,GMm^2}{\hbar}}
		\label{eq:cubic_explicit}
	\end{equation}
	
	%=============================================================
	\section{Complete Gravitational Wavefunction}
	%=============================================================
	
	\subsection{Normalized Form}
	
	\begin{equation}
		\boxed{\psi(x) = \frac{\alpha_G}{r}\begin{bmatrix}1\\0\\ p_zc/\Delta\\ (p_x+ip_y)c/\Delta\end{bmatrix}e^{-iS(r)/\hbar}}
		\label{eq:normalized_wavefunction}
	\end{equation}
	
	where $\alpha_G = e^{i\pi/4}\sqrt{c\hbar/(2GMm)}$,
	\begin{equation}
		\Delta = E + mc^2 + \int\rho c^2\,dV - \int P\,dV - \frac{J^2}{2mr^2} + V_{EM} + \rho_\Lambda,
	\end{equation}
	and $|p|$ satisfies Eq.~\eqref{eq:cubic_explicit}.
	
	\subsection{Hamilton-Jacobi Connection}
	
	The phase $S(r)$ satisfies:
	\begin{equation}
		\frac{1}{2m}\!\left(\frac{dS}{dr}\right)^2 + V(r) = E
	\end{equation}
	with $V(r)=-GMm/r$.  With $p(r)=dS/dr$:
	\begin{equation}
		F(r) = -\frac{dV}{dr} = -\frac{GMm}{r^2}.
	\end{equation}
	This recovers Newton's law of gravitation.
	
	%=============================================================
	\section{Relation to Gravitational Potential}
	%=============================================================
	
	For completeness, with $\Phi(r)=-GMm/r$:
	
	\subsection{Coupling Constant in Terms of Potential}
	
	\begin{equation}
		|\alpha_G|^2 = \frac{c\hbar}{2GMm} = -\frac{c\hbar}{2r\Phi(r)}
	\end{equation}
	
	\subsection{Wavefunction Amplitude}
	
	From the exact bridge $|\psi|^2 \cdot p \cdot r^2 = C$ with $p=\sqrt{-2m\Phi(r)}$ ($\Phi<0$) and $|\psi|^2 = |\alpha_G|^2(1+p^2c^2/\Delta^2)/r^2$:
	\begin{equation}
		|\psi|^2 \approx \frac{C}{r^2\,|p|} = \frac{mc/\sqrt{2}}{r^2\sqrt{2m|\Phi(r)|}} = \frac{c\sqrt{m}}{2r^2\sqrt{|\Phi(r)|}}.
	\end{equation}
	
	\begin{equation}
		\boxed{|\psi|^2 = \frac{c\sqrt{m}}{2\,r^2\sqrt{|\Phi(r)|}}}
	\end{equation}
	
	Or in proportional form:
	\begin{equation}
		\boxed{|\psi|^2 \propto \frac{1}{r^2\sqrt{|\Phi(r)|}}}
	\end{equation}
	
	The wavefunction probability density thus has two contributions: the $1/r^2$ geometric falloff from the spherical wave, and a $1/\sqrt{|\Phi(r)|}$ dependence on the gravitational potential.  Both reflect the $1/r^2$ inverse-square structure of three-dimensional gravity.
	
	%=============================================================
	\section{Stress-Energy and Einstein's Equations}
	%=============================================================
	
	\subsection{Coarse-Grained Stress-Energy}
	
	After averaging over scale $\ell$:
	\begin{equation}
		\langle T^{\mu\nu}\rangle = \rho c^2 u^\mu u^\nu + Pg^{\mu\nu} + \tau^{\mu\nu}_{\mathrm{spin}} + \tau^{\mu\nu}_{\mathrm{EM}} + \Lambda g^{\mu\nu}
		\label{eq:T_fluid}
	\end{equation}
	
	where $\rho=m\langle R^2\rangle$, $P$ from $P(\bar{\psi}\psi)$, $\tau^{\mu\nu}_{\mathrm{spin}}$ from $J_{\mu\nu}J^{\mu\nu}$, $\tau^{\mu\nu}_{\mathrm{EM}}$ from $F_{\mu\nu}F^{\mu\nu}$, and $\Lambda$ from $\rho_\Lambda$.
	
	Equation~\eqref{eq:T_fluid} has the structure of the standard Einstein stress-energy tensor.
	
	\subsection{Connection to General Relativity}
	
	In the weak-field limit, this sources the metric perturbation:
	\begin{equation}
		h_{\mu\nu} = \frac{16\pi G}{c^4}\int\frac{T_{\mu\nu}(x')}{|x-x'|}\,d^3x'
	\end{equation}
	
	The $g_{00}$ component gives:
	\begin{equation}
		g_{00} = -1 + \frac{2\Phi}{c^2}
	\end{equation}
	where $\Phi=-GM/r$ is the Newtonian potential per unit mass.
	
	%=============================================================
	\section{The Four-Component General Solution and the Gravitational Fine Structure Constant}
	\label{sec:four_component}
	%=============================================================
	
	\subsection{Complete Spinor Solution}
	
	\begin{equation}
		\psi = \frac{\alpha_G}{r}\Bigg[
		\psi_0\begin{bmatrix}1\\0\\ p_zc/\Delta\\ (p_x+ip_y)c/\Delta\end{bmatrix}e^{-iS/\hbar}
		+\psi_1\begin{bmatrix}0\\1\\ (p_x-ip_y)c/\Delta\\ -p_zc/\Delta\end{bmatrix}e^{-iS/\hbar}
		+\psi_2\begin{bmatrix}p_zc/\Delta'\\ (p_x+ip_y)c/\Delta'\\1\\0\end{bmatrix}e^{+iS/\hbar}
		+\psi_3\begin{bmatrix}(p_x-ip_y)c/\Delta'\\ -p_zc/\Delta'\\0\\1\end{bmatrix}e^{+iS/\hbar}
		\Bigg]
		\label{eq:four_component_wavefunction}
	\end{equation}
	
	where the universal normalization is:
	\begin{equation}
		\boxed{\alpha_G = e^{i\pi/4}\sqrt{\frac{c\hbar}{2GMm}}}
		\label{eq:universal_alphaG}
	\end{equation}
	
	This is the \textbf{gravitational fine structure constant}, in direct analogy to QED's $\alpha_{\rm em}=e^2/(4\pi\epsilon_0\hbar c)\approx 1/137$.
	
	The energy denominators are:
	\begin{align}
		\Delta  &= E + mc^2 + H \quad\text{(particle states)} \\
		\Delta' &= E - mc^2 - H \quad\text{(antiparticle states)}
	\end{align}
	where $H$ contains all stress-energy contributions from Eq.~\eqref{eq:hamiltonian}.
	
	\subsection{Physical Interpretation}
	
	The four components represent:
	\begin{itemize}
		\item $\psi_0,\psi_1$: Particle states (spin up/down), forward time evolution $e^{-iS/\hbar}$
		\item $\psi_2,\psi_3$: Antiparticle states (spin up/down), backward time evolution $e^{+iS/\hbar}$
	\end{itemize}
	
	All four states share the universal coupling $\alpha_G$, which satisfies:
	\begin{equation}
		[\alpha_G^2] = \left[\frac{c\hbar}{GMm}\right] = \frac{(LT^{-1})(ML^2T^{-1})}{(L^3M^{-1}T^{-2})\cdot M\cdot M} = 1 \quad\checkmark
	\end{equation}

	\begin{tcolorbox}[enhanced,colback=yellow!8,colframe=black,
	  title={\textbf{Note on the dimensions of $\alpha_G$ and $\psi$}},
	  fonttitle=\bfseries,breakable]
	The check $[\alpha_G^2] = 1$ establishes that $\alpha_G$ is \emph{dimensionless}.
	The wavefunction is written $\psi = (\alpha_G/r)\,u\,e^{-iS/\hbar}$, giving
	$[\psi] = L^{-1}[\,u\,]$.  For $|\psi|^2$ to serve as a probability density in 3D,
	one requires $[\psi] = L^{-3/2}$, which would demand $[\,u\,] = L^{-1/2}$.  In the
	normalisation scheme of this paper, $u$ is an $\mathcal{O}(1)$ dimensionless spinor
	envelope, so $[\psi] = L^{-1}$ and $|\psi|^2$ has units $L^{-2}$.

	This is consistent: the probability flux $\Phi = 4\pi r^2 \cdot j^r = 4\pi|\alpha_G|^2 p$
	is finite and conserved, and the current-conservation bridge
	$|\psi|^2 \cdot p \cdot r^2 = C$ is dimensionally homogeneous with $[C] = \mathrm{momentum}$.
	The factor $r^2$ in the bridge is the spherical area element, not a normalisation
	artefact.  Throughout this paper $|\psi|^2$ is used as a \emph{flux density}
	(probability current amplitude) rather than a volumetric probability density;
	quantities such as $|\psi|^2 \cdot p$ have the correct dimensions for a momentum flux.
	Readers who prefer standard $L^{-3/2}$ normalisation should regard $u$ as carrying
	dimension $L^{-1/2}$ and $\alpha_G$ as dimensionless, with the $1/r$ factor supplying
	the remaining $L^{-1}$.
	\end{tcolorbox}
	
	%=============================================================
	\section{Probability Density: Full Four-Component Case}
	%=============================================================
	
	\subsection{Spinor Norms}
	
	With $\psi = (\alpha_G/r)\,u\,e^{iS/\hbar}$, the probability density is $|\psi|^2 = (|\alpha_G|^2/r^2)\,|u|^2$.  The spinor norms $|u|^2$ are:
	
	\textbf{Particle state $\psi_0$:}
	\begin{align}
		\left|\begin{bmatrix}1\\0\\ p_zc/\Delta\\ (p_x+ip_y)c/\Delta\end{bmatrix}\right|^2
		= 1 + \frac{|p|^2 c^2}{\Delta^2}
		\label{eq:norm_psi0}
	\end{align}
	
	\textbf{Particle state $\psi_1$:}
	\begin{equation}
		\left|\begin{bmatrix}0\\1\\ (p_x-ip_y)c/\Delta\\ -p_zc/\Delta\end{bmatrix}\right|^2 = 1 + \frac{|p|^2 c^2}{\Delta^2}
		\label{eq:norm_psi1}
	\end{equation}
	
	\textbf{Antiparticle state $\psi_2$:}
	\begin{equation}
		\left|\begin{bmatrix}p_zc/\Delta'\\ (p_x+ip_y)c/\Delta'\\1\\0\end{bmatrix}\right|^2 = 1 + \frac{|p|^2 c^2}{\Delta'^2}
		\label{eq:norm_psi2}
	\end{equation}
	
	\textbf{Antiparticle state $\psi_3$:}
	\begin{equation}
		\left|\begin{bmatrix}(p_x-ip_y)c/\Delta'\\ -p_zc/\Delta'\\0\\1\end{bmatrix}\right|^2 = 1 + \frac{|p|^2 c^2}{\Delta'^2}
		\label{eq:norm_psi3}
	\end{equation}
	
	\subsection{Time-Averaged Probability Density}
	
	Interference terms $\propto e^{\pm 2iS/\hbar}$ vanish in the semiclassical WKB limit.  With occupation coefficients:
	\begin{align}
		A &\equiv |\psi_0|^2+|\psi_1|^2 \quad\text{(particle occupation)} \\
		B &\equiv |\psi_2|^2+|\psi_3|^2 \quad\text{(antiparticle occupation)}
	\end{align}
	
	The full probability density with the spherical-wave $1/r$ factor is:
	\begin{equation}
		\boxed{
			|\psi|^2 = \frac{|\alpha_G|^2}{r^2}\!\left[A\!\left(1+\frac{|p|^2 c^2}{\Delta^2}\right)+B\!\left(1+\frac{|p|^2 c^2}{\Delta'^2}\right)\right]
		}
		\label{eq:prob_density_general}
	\end{equation}
	
	%=============================================================
	\section{Current Conservation Constraint and the Cubic Equation}
	%=============================================================
	
	From probability current conservation in spherical symmetry:
	\begin{equation}
		\boxed{|\psi|^2 = \frac{C}{|p|}}
	\end{equation}
	
	Equating with Eq.~\eqref{eq:prob_density_general} and multiplying by $|p|$:
	\begin{equation}
		|\alpha_G|^2\!\left[(A+B)|p| + c^2|p|^3\!\left(\frac{A}{\Delta^2}+\frac{B}{\Delta'^2}\right)\right] = C
	\end{equation}
	
	This yields a depressed cubic:
	\begin{equation}
		\boxed{\alpha|p|^3 + \beta|p| - C = 0}
		\label{eq:cubic_momentum}
	\end{equation}
	
	with:
	\begin{align}
		\alpha &= c^2|\alpha_G|^2\!\left(\frac{A}{\Delta^2}+\frac{B}{\Delta'^2}\right) = \frac{c^3\hbar}{2GMm}\!\left(\frac{A}{\Delta^2}+\frac{B}{\Delta'^2}\right) \label{eq:alpha_coeff}\\
		\beta  &= |\alpha_G|^2(A+B) = \frac{c\hbar}{2GMm}(A+B) \label{eq:beta_coeff}
	\end{align}
	
	Solution via Cardano's formula:
	\begin{equation}
		|p| = \left[\frac{C}{2\alpha}+\sqrt{\left(\frac{C}{2\alpha}\right)^2+\left(\frac{\beta}{3\alpha}\right)^3}\right]^{1/3}
		+ \left[\frac{C}{2\alpha}-\sqrt{\left(\frac{C}{2\alpha}\right)^2+\left(\frac{\beta}{3\alpha}\right)^3}\right]^{1/3}
		\label{eq:cardano_solution}
	\end{equation}
	
	%=============================================================
	\section{The Fundamental Momentum Scale $\mathcal{P}$}
	%=============================================================
	
	\subsection{Definition}
	
	For the single-component case ($A=1$, $B=0$):
	\begin{equation}
		\mathcal{P} \equiv \frac{C}{|\alpha_G|^2} = \frac{mc/\sqrt{2}}{c\hbar/(2GMm)} = \frac{\sqrt{2}\,GMm^2}{\hbar}.
	\end{equation}
	
	\begin{equation}
		\boxed{\mathcal{P} = \frac{\sqrt{2}\,GMm^2}{\hbar}}
		\label{eq:B_final}
	\end{equation}
	
	\subsection{Dimensional Analysis}
	
	\begin{equation}
		[\mathcal{P}] = \frac{[G][M][m^2]}{[\hbar]} = \frac{(L^3M^{-1}T^{-2})\cdot M\cdot M^2}{ML^2T^{-1}} = LMT^{-1} \quad\checkmark
	\end{equation}
	
	\subsection{Factorization}
	
	\begin{equation}
		\boxed{C = |\alpha_G|^2\,\mathcal{P}}
		\label{eq:C_factorization}
	\end{equation}
	
	The probability flux equals the fundamental momentum scale suppressed by the squared modulus of the gravitational coupling.
	
	\subsection{The Gravitational Bohr Momentum}
	
	$\mathcal{P}$ is the gravitational analogue of the Bohr momentum.  In atomic physics:
	\begin{equation}
		p_{\mathrm{Bohr}} = \frac{me^2}{4\pi\epsilon_0\hbar} \sim \alpha_{\mathrm{em}}\,mc.
	\end{equation}
	
	Similarly:
	\begin{equation}
		\mathcal{P} = \frac{\sqrt{2}\,GMm^2}{\hbar} = \frac{m}{\hbar}\cdot\frac{\sqrt{2}\,GMm}{\hbar} \cdot \hbar
		= \frac{\sqrt{2}\,m\,E_{\mathrm{grav}}}{\hbar c}\cdot c
	\end{equation}
	
	\subsection{Regime Classification}
	
	\paragraph{Gravity-dominated:} $(C/2\alpha)^2\gg(\beta/3\alpha)^3$:
	\begin{equation}
		|p| \approx (C/\alpha)^{1/3} \sim \mathcal{P}^{1/3}.
	\end{equation}
	
	\paragraph{Linear-dominated:} $(\beta/3\alpha)^3\gg(C/2\alpha)^2$:
	\begin{equation}
		|p| \approx \frac{C}{\beta} = \frac{\mathcal{P}}{A+B}.
	\end{equation}
	
	%=============================================================
	\section{Antiparticle Suppression and Hawking Radiation}
	%=============================================================
	
	Near a black hole horizon at Hawking temperature $T_H=\hbar c^3/(8\pi GMk_B)$, the antiparticle occupation is:
	\begin{equation}
		\frac{B}{A} \sim \exp\!\left(-\frac{2mc^2}{k_BT_H}\right) = \exp\!\left(-\frac{16\pi GMm}{\hbar c}\right) = \exp\!\left(-\frac{8\pi}{|\alpha_G|^2}\right).
	\end{equation}
	
	For $|\alpha_G|^2\gg1$ (weak gravity): $B\approx0$, and the theory reduces to the single-component particle sector.\\
	For $|\alpha_G|^2\sim1$ (Planck scale): antiparticle occupation becomes of order unity, which is consistent with the Hawking radiation regime.
	
	%=============================================================
	\section{The Graviton Scalar and Special Relativistic Structure}
	\label{sec:graviton_scalar}
	%=============================================================
	
	\subsection{Motivation from Spinor Components}
	
	The lower spinor components have the form $p_ic/\Delta$---dimensionless, analogous to $\beta=v/c$ in special relativity.
	
	\subsection{Definition of the Graviton Scalar}
	
	\begin{equation}
		\boxed{\sigma_i \equiv \frac{p_i c}{\Delta}}
		\label{eq:graviton_scalar}
	\end{equation}
	
	For particle states $\Delta=E+mc^2+H$; for antiparticle states $\Delta'=E-mc^2-H$.
	
	\subsection{Physical Interpretation}
	
	\begin{enumerate}
		\item \textbf{Dimensionless rapidity parameter:} $\sigma$ measures how quantum-gravitational a configuration is.
		\item \textbf{Reduction to SR:} In the non-relativistic limit $\Delta\approx mc^2$, $p=mv$:
		\begin{equation}
			\sigma = \frac{pc}{mc^2} = \frac{v}{c} = \beta.
		\end{equation}
		\item \textbf{Regime classification:}
		\begin{itemize}
			\item $\sigma\ll1$: Classical/Newtonian gravity
			\item $\sigma\sim1$: Full quantum gravity required
			\item $\sigma\gg1$: Strongly quantum gravitational
		\end{itemize}
	\end{enumerate}
	
	\subsection{Four-Component Solution in $\sigma$-Basis}
	
	\begin{equation}
		\psi = \frac{\alpha_G}{r}\Bigg[
		\psi_0\begin{bmatrix}1\\0\\\sigma_z\\\sigma_x+i\sigma_y\end{bmatrix}e^{-iS/\hbar}
		+\psi_1\begin{bmatrix}0\\1\\\sigma_x-i\sigma_y\\-\sigma_z\end{bmatrix}e^{-iS/\hbar}
		+\psi_2\begin{bmatrix}\sigma_z'\\\sigma_x'+i\sigma_y'\\1\\0\end{bmatrix}e^{+iS/\hbar}
		+\psi_3\begin{bmatrix}\sigma_x'-i\sigma_y'\\-\sigma_z'\\0\\1\end{bmatrix}e^{+iS/\hbar}
		\Bigg]
	\end{equation}
	
	\subsection{Gravitational Lorentz Factor}
	
	\begin{equation}
		\boxed{|u|^2 = 1+|\boldsymbol{\sigma}|^2 = 1+\frac{p^2c^2}{\Delta^2}}
		\label{eq:gravitational_lorentz_factor}
	\end{equation}
	
	\begin{center}
		\begin{tabular}{c|c}
			\textbf{Special Relativity} & \textbf{Quantum Gravity} \\
			\hline
			$\beta = v/c$ & $\sigma = pc/\Delta$ \\
			$\gamma^2 = 1/(1-\beta^2)$ & $|u|^2 = 1+\sigma^2$ \\
			Diverges at $v\to c$ & Grows as $p\to\Delta$ \\
			Time dilation & Wavefunction quantum effects \\
		\end{tabular}
	\end{center}
	
	\subsection{Cubic Equation in $\sigma$-Form}
	
	Multiplying $|p|+c^2|p|^3/\Delta^2=\mathcal{P}$ by $c/\Delta$ and using $\sigma=pc/\Delta$:
	\begin{equation}
		\boxed{\sigma + \sigma^3 = \frac{\mathcal{P}c}{\Delta}}
		\label{eq:cubic_sigma}
	\end{equation}
	
	This is a nonlinear dispersion relation; within the present framework, the gravitational dynamics are nonlinear at the level of the wavefunction phase.
	
	\subsection{Post-Newtonian Expansion}
	
	Expanding in powers of $\sigma$:
	\begin{equation}
		|\psi|^2 = \frac{|\alpha_G|^2}{r^2}(1+\sigma^2+\mathcal{O}(\sigma^4))
	\end{equation}
	
	Each power corresponds to a post-Newtonian order:
	\begin{itemize}
		\item $\sigma^0$: Newtonian gravity
		\item $\sigma^2$: 1PN corrections
		\item $\sigma^3$: Nonlinear quantum corrections (from cubic equation)
	\end{itemize}
	
	%=============================================================
	\section{Graviton Frequency}
	%=============================================================
	
	From $\mathcal{P}=\sqrt{2}\,GMm^2/\hbar$, the graviton frequency is:
	\begin{equation}
		\omega_g = \frac{\mathcal{P}}{\hbar} = \frac{\sqrt{2}\,GMm^2}{\hbar^2}
	\end{equation}
	
	\begin{equation}
		\boxed{\mathcal{P} = \hbar\omega_g}
	\end{equation}
	
	%=============================================================
	\section{Field Energy and Momentum}
	%=============================================================
	
	\subsection{Length Scales}
	
	Define the Schwarzschild radius and reduced Compton wavelength:
	\begin{align}
		R_s &\equiv \frac{2GM}{c^2}, \qquad \lambda_c \equiv \frac{\hbar}{mc}.
	\end{align}
	
	Key relation:
	\begin{equation}
		\frac{R_s}{\lambda_c} = \frac{2GMm}{\hbar c} = \frac{1}{|\alpha_G|^2}
		\label{eq:Rs_lambda_relation}
	\end{equation}
	since $|\alpha_G|^2 = c\hbar/(2GMm)$.
	
	\subsection{Field Energy}
	
	For massless gravitons ($E=pc$):
	\begin{equation}
		E_{\mathrm{field}} = \mathcal{P}\cdot c = \frac{\sqrt{2}\,GMm^2c}{\hbar}.
	\end{equation}
	
	Substituting $GM=R_sc^2/2$ and $\hbar=mc\lambda_c$:
	\begin{equation}
		\boxed{E_{\mathrm{field}} = \frac{1}{\sqrt{2}}\,mc^2\left(\frac{R_s}{\lambda_c}\right) = \frac{mc^2}{\sqrt{2}\,|\alpha_G|^2}}
		\label{eq:E_field}
	\end{equation}
	
	\subsection{Field Momentum}
	
	\begin{equation}
		\boxed{P_{\mathrm{field}} = \mathcal{P} = \frac{\sqrt{2}\,GMm^2}{\hbar} = \frac{1}{\sqrt{2}}\,mc\left(\frac{R_s}{\lambda_c}\right) = \frac{mc}{\sqrt{2}\,|\alpha_G|^2}}
		\label{eq:P_field}
	\end{equation}
	
	\subsection{Physical Interpretation}
	
	Equation~\eqref{eq:E_field} connects three characteristic scales:
	\begin{itemize}
		\item \textbf{Quantum mechanics}: $\lambda_c=\hbar/(mc)$
		\item \textbf{Special relativity}: $mc^2$ (rest energy)
		\item \textbf{General relativity}: $R_s=2GM/c^2$
	\end{itemize}
	The ratio $R_s/\lambda_c = 1/|\alpha_G|^2$ is the single dimensionless parameter controlling the relative strength of gravitational and quantum effects.
	
	\subsection{Regime Classification}
	
	\begin{enumerate}
		\item \textbf{Quantum regime} ($|\alpha_G|^2\gg1$, i.e.\ $c\hbar\gg2GMm$):\\
		Field energy $\ll mc^2$; gravity negligible.
		\item \textbf{Quantum gravity regime} ($|\alpha_G|^2\sim1$):\\
		Field energy $\sim mc^2$; full four-component structure necessary.
		\item \textbf{Classical regime} ($|\alpha_G|^2\ll1$, i.e.\ $2GMm\gg c\hbar$):\\
		Field energy $\gg mc^2$; antiparticles exponentially suppressed.
	\end{enumerate}
	
	\subsection{Numerical Examples}
	
	\textbf{Earth--electron system:}
	\begin{align}
		|\alpha_G|^2 &= \frac{(3\times10^8)(1.055\times10^{-34})}{2(6.67\times10^{-11})(5.97\times10^{24})(9.11\times10^{-31})} \approx 4.35\times10^{-10}.
	\end{align}
	Classical GR regime ($|\alpha_G|^2\ll1$): gravitational field energy $\gg mc^2$.
	
	\textbf{Planck-mass particle ($M=m=M_{\rm Pl}$):}
	\begin{equation}
		|\alpha_G|^2 = \frac{c\hbar}{2GM_{\rm Pl}^2} = \frac{c\hbar}{2G(\hbar c/G)} = \frac{1}{2}.
	\end{equation}
	Quantum gravity regime: $E_{\mathrm{field}} = \sqrt{2}\,M_{\rm Pl}c^2\approx 1.73\times10^{19}$~GeV.
	
	\subsection{Connection to $\sigma$-Field Formalism}
	
	The covariant vector field $\sigma_\mu(x)=\partial_\mu S(x)/(mc)$ satisfies:
	\begin{equation}
		|\sigma| \sim \frac{|p|}{mc} \sim \frac{\mathcal{P}}{mc} = \frac{\sqrt{2}\,GMm}{c\hbar} = \frac{\sqrt{2}}{|\alpha_G|^2}\cdot\frac{1}{2} = \frac{1}{\sqrt{2}\,|\alpha_G|^2}.
	\end{equation}
	
	%=============================================================
	\section{The Exact Quantum Force Law and Higher-Order Corrections}
	\label{sec:UV_identification}
	%=============================================================

	\subsection{The exact phase-field formula}

	The momentum and energy fields in the weak-field limit satisfy:
	\begin{equation}
		E(r) = \frac{mc^2}{\alpha_G^2\,e^{2imcr/\hbar}}.
		\label{eq:E_phase}
	\end{equation}
	Substituting $\alpha_G^2 = ic\hbar/(2GMm)$:
	\begin{equation}
		\boxed{E(r) = -\frac{GMm}{e^{2imcr/\hbar}} = -GMm\,e^{-2imcr/\hbar}.}
		\label{eq:E_exact_exp}
	\end{equation}
	Factoring $e^{it} = it\cdot(\sin t/t)$ with $t = 2r/\lambda_C = 2mcr/\hbar$:
	\begin{equation}
		e^{it} = \frac{2imcr}{\hbar}\cdot\frac{\sin(2r/\lambda_C)}{2r/\lambda_C},
	\end{equation}
	so the exact energy can be written:
	\begin{equation}
		\boxed{E(r) = -\frac{GMm}{r}\cdot\frac{2r/\lambda_C}{\sin(2r/\lambda_C)}.}
		\label{eq:E_sinh}
	\end{equation}
	Note: $E(r)$ in~\eqref{eq:E_sinh} is \emph{real} and recovers $-GMm/r$
	as $r\to0$ (since $x/\sin x\to1$). It is the exact gravitational potential
	at the quantum-gravity scale, valid for $r < \pi\lambda_C/2$.

	\subsection{Derivation of Newton's law}

	At leading order $r\ll\lambda_C$, using $\sin(2r/\lambda_C)\approx 2r/\lambda_C$:
	\begin{equation}
		E(r) \approx -\frac{GMm}{r}, \qquad F = -\frac{dE}{dr} = \frac{GMm}{r^2}.\quad\checkmark
	\end{equation}

	\subsection{Higher-order corrections: term-by-term and exact}

	\paragraph{Term-by-term from the Taylor expansion of $e^{it}$.}
	Expanding $e^{it} = it(1 - t^2/6 + t^4/120 - \cdots)$:
	\begin{equation}
		E(r) = -\frac{GMm}{r}\cdot\frac{t}{\sin t}
		     = -\frac{GMm}{r}\!\left[1 + \frac{t^2}{6} + \frac{7t^4}{360}
		       + \frac{31t^6}{15120} + \frac{127t^8}{604800} + \cdots\right]
	\end{equation}
	where $t = 2r/\lambda_C$.  Taking $F=-dE/dr$ and using $dt/dr = 2/\lambda_C$:
	\begin{equation}
		F(r) = \frac{GMm}{r^2}\!\left[1 + \frac{9}{2}\!\left(\frac{r}{\lambda_C}\right)^{\!2}
		      + \cdots\right].
	\end{equation}
	More precisely, each coefficient $a_k$ in
	$F = (GMm/r^2)\sum_{k=0}^\infty a_k t^{2k}$
	is extracted from the exact closed form below.

	\begin{tcolorbox}[enhanced,breakable,
	  colback=white,colframe=black,
	  title={\textbf{Note on the force law expansion coefficients}},
	  fonttitle=\bfseries]
	The summary formula $F(r) = (GMm/r^2)[1 + (9/2)(\lambda_C/r)^2 + \cdots]$ in
	Section~\ref{sec:summary_force} is expressed in powers of $(\lambda_C/r)$, i.e.\
	corrections grow as the distance \emph{decreases} toward $\lambda_C$.  In contrast,
	the series here is in powers of $t^{2k} = (2r/\lambda_C)^{2k}$, which grows with
	\emph{increasing} $r$.  The leading quantum correction term from differentiating
	$E = -(GMm/r)(t/\sin t)$ is:
	\begin{equation}
	  F(r) = \frac{GMm}{r^2}\left[1 + \frac{9}{2}\left(\frac{r}{\lambda_C}\right)^2
	         + \mathcal{O}\!\left(\frac{r^4}{\lambda_C^4}\right)\right],
	         \quad t = \frac{2r}{\lambda_C} \ll 1.
	\end{equation}
	This expansion is valid only for $r \ll \lambda_C$ (the deep quantum regime,
	$t \ll 1$).  The coefficient $+9/2$ multiplies $(r/\lambda_C)^2$, not
	$(\lambda_C/r)^2$, and is the leading sub-leading correction to Newton's law
	within the quantum-gravity regime.  Outside this domain (i.e.\ for
	$r \gg \lambda_C$, ordinary macroscopic gravity), neither this series nor the
	exact formula~\eqref{eq:F_exact_closed} applies: Newton's law is instead
	recovered from the current-conservation bridge and the QGD field equation
	(see Remark~\ref{rem:regimes}).
	\end{tcolorbox}

	\paragraph{Exact closed form.}
	From~\eqref{eq:E_sinh}, differentiating exactly:
	\begin{equation}
		\boxed{F(r) = \frac{GMm}{r^2}\cdot\frac{t^2\cos t}{\sin^2 t},
		\qquad t = \frac{2r}{\lambda_C}.}
		\label{eq:F_exact_closed}
	\end{equation}
	Equation~\eqref{eq:F_exact_closed} is the force law within the quantum-gravity regime ($r < \pi\lambda_C/2$) of QGD.

	\subsection{Series coefficients: a closed formula for all orders}

	Expanding $t^2\cos t/\sin^2 t$ in powers of $t$ (equivalent to powers of
	$\lambda_C/r$ via $t = 2r/\lambda_C$) using the known series
	$t^2\cot^2(t/2) = \sum_{n=0}^\infty c_n t^{2n}$:
	\begin{equation}
		F(r) = \frac{GMm}{r^2}\!\left[
		  1
		  - \frac{t^2}{6}
		  - \frac{7t^4}{120}
		  - \frac{31t^6}{3024}
		  - \frac{127t^8}{86400}
		  - \frac{73t^{10}}{380160}
		  - \cdots
		\right].
	\end{equation}

	The coefficients of $t^{2k}$ are related to the \textbf{Bernoulli numbers}
	$B_{2k}$ by the identity for $t/\sin t$:
	\begin{equation}
		\frac{t}{\sin t} = \sum_{k=0}^{\infty}
		  \frac{(2-2^{2k})\,B_{2k}}{(2k)!}\,(-1)^{k+1}\,t^{2k}
		= 1 + \frac{t^2}{6} + \frac{7t^4}{360} + \frac{31t^6}{15120} + \cdots
	\end{equation}
	so that the force coefficients $a_k$ in $F = (GMm/r^2)\sum_{k\geq0}a_k t^{2k}$
	satisfy the closed formula:
	\begin{equation}
		\boxed{a_k = \frac{d}{dt^{2k}}\!\left[\frac{t^2\cos t}{\sin^2 t}\right]_{t=0}
		     = 1 - (2k+1)\,\frac{(2-2^{2k})(-1)^{k+1}B_{2k}}{(2k)!}}
		\label{eq:a_k_closed}
	\end{equation}
	(with $a_0=1$, $B_0=1$, $B_2=1/6$, $B_4=-1/30$, $B_6=1/42$, \ldots).

	\begin{table}[h]
	\centering
	\renewcommand{\arraystretch}{1.3}
	\caption{First six quantum-gravity force coefficients from~\eqref{eq:F_exact_closed}.
	  These are coefficients of $t^{2k} = (2r/\lambda_C)^{2k}$ in the expansion of
	  $F(r)/(GMm/r^2)$, valid for $r \ll \lambda_C$ (quantum regime, $t\ll1$).
	  The exact formula~\eqref{eq:F_exact_closed} must be used for $r \sim \lambda_C$.}
	\begin{tabular}{@{}ccccc@{}}
	\toprule
	$k$ & Coefficient of $t^{2k}$ & $a_k$ & In $(r/\lambda_C)^{2k}$ & In $(\lambda_C/r)^{2k}$ \\
	\midrule
	0 & $t^0$  & $1$         & 1  & 1 \\
	1 & $t^2$  & $-1/6$      & $-(2/3)(r/\lambda_C)^2$ & Leading quantum term \\
	2 & $t^4$  & $-7/120$    & ...   &   \\
	3 & $t^6$  & $-31/3024$  & ...   &   \\
	4 & $t^8$  & $-127/86400$& ...   &   \\
	5 & $t^{10}$&$-73/380160$& ...   &   \\
	\bottomrule
	\end{tabular}
	\end{table}

	\subsection{Domain and physical interpretation}

	\begin{itemize}
		\item \textbf{Valid range:} $r < \pi\lambda_C/2 \approx 1.57\,\lambda_C$.
		  At $r = \pi\lambda_C/2$: $\sin(2r/\lambda_C) = 0$, force diverges ---
		  a \emph{quantum gravity horizon} at the Compton scale.
		\item \textbf{Classical limit} $r\gg\lambda_C$: this formula does \emph{not}
		  apply.  Newton's law in the classical regime is recovered from the
		  current-conservation bridge and the QGD field equation (see
		  Remark~\ref{rem:regimes}).
		\item \textbf{Quantum regime} $r\sim\lambda_C$: all orders are comparable;
		  the exact formula~\eqref{eq:F_exact_closed} must be used.
		\item \textbf{Crossover scale:} quantum correction equals Newtonian
		  at $r_c = (3/\sqrt{2})\lambda_C \approx 2.12\,\lambda_C$. For an
		  electron: $r_c \approx 8\times10^{-13}$\,m.
	\end{itemize}

	\begin{table}[h]
		\centering
		\begin{tabular}{@{}lcc@{}}
			\toprule
			Regime & Distance & Dominant term \\
			\midrule
			Classical   & $r\gg\lambda_C$ & Use QGD field equation $\square_g\sigma=0$ \\
			Transition  & $r\sim2\lambda_C$ & Exact formula~\eqref{eq:F_exact_closed} \\
			Quantum     & $r\lesssim\lambda_C$ & All orders of series; diverges at $r=\pi\lambda_C/2$ \\
			\bottomrule
		\end{tabular}
	\end{table}
	
	%=============================================================
	\section{Mapping Table: Action $\to$ Stress-Energy $\to$ Wavefunction}
	%=============================================================
	
	\begin{center}
		\begin{tabular}{|l|l|l|l|}
			\hline
			\textbf{Action Term} & \textbf{Stress-Energy} & \textbf{Enters $\Delta$ as} & \textbf{Effect} \\
			\hline
			$-m\bar{\psi}\psi$ & $\rho c^2 u^\mu u^\nu$ & $+mc^2$ & Test mass \\
			$\frac{i}{2}\bar{\psi}\gamma^\mu\overleftrightarrow{\partial}_\mu\psi$ & $\rho c^2 u^\mu u^\nu$ & $+\int\rho c^2\,dV$ & Source mass \\
			$-P(\bar{\psi}\psi)$ & $Pg^{\mu\nu}$ & $-\int P\,dV$ & Pressure work \\
			$-\frac{1}{2M^2}J_{\mu\nu}J^{\mu\nu}\bar{\psi}\psi$ & $\tau^{\mu\nu}_{\rm spin}$ & $-J^2/(2mr^2)$ & Frame-dragging \\
			$-\frac{1}{4}F_{\mu\nu}F^{\mu\nu}$ & $\tau^{\mu\nu}_{\rm EM}$ & $+V_{EM}$, $p\to p-eA/c$ & EM coupling \\
			$-\rho_\Lambda$ & $\rho_\Lambda g^{\mu\nu}$ & $+\rho_\Lambda$ & Vacuum energy \\
			\hline
		\end{tabular}
	\end{center}
	
	%=============================================================
	\section{Summary of Key Results}
	\label{sec:summary_force}
	%=============================================================
	
	\subsection{The Gravitational Fine Structure Constant}
	
	\begin{equation}
		\boxed{\alpha_G = e^{i\pi/4}\sqrt{\frac{c\hbar}{2GMm}}, \qquad \alpha_G^2 = \frac{ic\hbar}{2GMm}, \qquad |\alpha_G|^2 = \frac{c\hbar}{2GMm}}
	\end{equation}
	
	\emph{When a real value is required:} use $|\alpha_G|^2$ for densities/fluxes, or $\mathrm{Re}(\alpha_G)=\sqrt{c\hbar/(4GMm)}$ for amplitudes.
	
	\subsection{Derived Scales}
	
	\begin{align}
		\mathcal{P} &= \frac{\sqrt{2}\,GMm^2}{\hbar} \quad\text{(fundamental momentum scale)} \\
		C &= \frac{mc}{\sqrt{2}} = |\alpha_G|^2\mathcal{P} \quad\text{(probability flux)} \\
		\omega_g &= \frac{\mathcal{P}}{\hbar} = \frac{\sqrt{2}\,GMm^2}{\hbar^2} \quad\text{(graviton frequency)} \\
		\frac{R_s}{\lambda_c} &= \frac{1}{|\alpha_G|^2} \quad\text{(length-scale ratio)}
	\end{align}
	
	\subsection{The Cubic Momentum Equation}
	
	\begin{equation}
		\boxed{|p| + \frac{c^2|p|^3}{\Delta^2} = \mathcal{P}}
	\end{equation}
	
	Interpolates between:
	\begin{itemize}
		\item Linear regime ($|p|\ll\Delta$): $|p|\approx\mathcal{P}$ (classical)
		\item Cubic regime ($|p|\sim\Delta/c$): $c^2|p|^3/\Delta^2\approx\mathcal{P}$ (quantum)
	\end{itemize}
	
	\subsection{Field Energy and Momentum}
	
	\begin{align}
		E_{\mathrm{field}} &= \frac{1}{\sqrt{2}}\,mc^2\left(\frac{R_s}{\lambda_c}\right) = \frac{mc^2}{\sqrt{2}\,|\alpha_G|^2} \\
		P_{\mathrm{field}} &= \frac{1}{\sqrt{2}}\,mc\left(\frac{R_s}{\lambda_c}\right) = \frac{mc}{\sqrt{2}\,|\alpha_G|^2}
	\end{align}
	
	\subsection{Force Law}
	\label{sec:summary_force_law}

	The exact quantum-gravity force law, valid for $r < \pi\lambda_C/2$ (quantum regime only):
	\begin{equation}
		F(r) = \frac{GMm}{r^2}\cdot\frac{t^2\cos t}{\sin^2 t}, \qquad t = \frac{2r}{\lambda_C}.
	\end{equation}
	For $r \ll \lambda_C$ the Taylor expansion gives:
	\begin{equation}
		F(r) \approx \frac{GMm}{r^2}\left[1 - \frac{1}{6}t^2 + \mathcal{O}(t^4)\right]
		           = \frac{GMm}{r^2}\left[1 - \frac{2}{3}\!\left(\frac{r}{\lambda_C}\right)^{\!2} + \mathcal{O}(r^4)\right],
		\label{eq:F_leading_corrected}
	\end{equation}
	where $t = 2r/\lambda_C$.  The leading quantum correction is \emph{negative}: gravity
	is slightly \emph{weaker} than Newton within the quantum domain $r \lesssim \lambda_C$.
	The correction coefficient $-1/6$ in $t^2$ (equivalently $-2/3$ in $(r/\lambda_C)^2$)
	is confirmed by direct symbolic differentiation of $E(r) = -(GMm/r)(t/\sin t)$
	(verified with SymPy; see Table in Section~\ref{sec:UV_identification}).
	The earlier expression $(9/2)(\lambda_C/r)^2$ appearing in some drafts is incorrect
	and has been removed.
	For macroscopic $r \gg \lambda_C$, Newton's law $F = GMm/r^2$ is recovered via the QGD field equation (see Remark~\ref{rem:regimes}).

	\subsection{The Gravitational Wavefunction}
		
		In the QGD framework, gravitational phenomena are proposed to emerge from macroscopic quantum coherence. The gravitational wavefunction takes the form:
		\begin{equation}\label{eq:wavefunction}
			\psi(x^\mu) = R(x)e^{-\frac{i}{\hbar}S(x)} \cdot u
		\end{equation}
		
		where $S(x)$ is the classical action, $R(x)$ is amplitude, $u$ is a spinor.
		
		\subsection{The Fundamental Quantum Constraint}
		
		Wavefunction-Field Relation]
			The probability density and gravitational field strength satisfy:
			\begin{equation}\label{eq:quantum_constraint}
				\boxed{|\psi(x)|^2 \sigma_\mu x^\mu = \frac{J}{\hbar}}
			\end{equation}
			where $J$ is a conserved quantum number, $\sigma_\mu = \nabla_\mu S/c$, and $x^\mu$ is the spacetime position.
		
			From probability current conservation $\nabla \cdot (|\psi|^2\nabla S) = 0$ in spherical symmetry:
			\begin{equation}
				|\psi|^2 = \frac{C}{|\mathbf{p}|} = \frac{C\lambda}{\hbar}
			\end{equation}
			
			Using $\sigma_r = r/\lambda$ (natural scalar interpretation) and $p = \hbar/\lambda$:
			\begin{equation}
				|\psi|^2 = \frac{C}{r\hbar\sigma_r}
			\end{equation}
			
			Multiplying by $\sigma_r r$ and generalizing covariantly yields (\ref{eq:quantum_constraint}).
		
		\textbf{Physical meaning:}
		\begin{itemize}
			\item Left side: (Quantum probability) $\times$ (Number of wavelengths from origin)
			\item Right side: Quantized phase space volume
			\item Strong field ($\sigma$ large) $\Rightarrow$ suppressed amplitude
			\item Weak field ($\sigma$ small) $\Rightarrow$ enhanced amplitude
		\end{itemize}
		
		Equation~\eqref{eq:quantum_constraint} may be interpreted as a quantum-gravitational phase-space constraint: the left side is the product of probability density and a measure of distance in units of the de~Broglie wavelength, and the right side is a conserved quantum number.
		
		Connection to sigma Field
		
		The phase gradient defines the gravitational field:
		\begin{equation}
			\boxed{\sigma_\mu \equiv \frac{1}{c}\partial_\mu S = \frac{p_\mu}{c}}
		\end{equation}
		
		$\sigma_\mu$ is the \textbf{gravitational phase field}, carrying momentum $p_\mu = \hbar\partial_\mu S/\hbar = \hbar\sigma_\mu/c$.


	%=============================================================
	\section{The Fundamental Action of QGD}
	\label{sec:qgd_action}
	%=============================================================
	
	The preceding sections derived the graviton field $\sigma_\mu$ from the Dirac wavefunction and showed that every stress-energy source enters through the unified denominator $\Delta$.  We now state the action principle that governs $\sigma_\mu$ as a dynamical field.
	
	Within QGD the metric is treated as a composite field rather than a fundamental degree of freedom:
	\begin{equation}
		g_{\mu\nu}(\sigma) = T^\alpha{}_\mu\,T^\beta{}_\nu\!\left(M_{\alpha\beta}\circ\left[\eta_{\alpha\beta} - \sum_a\varepsilon_a\,\sigma_\alpha^{(a)}\sigma_\beta^{(a)} - \kappa\,\ell_Q^2\,\partial_\alpha\sigma^\gamma\partial_\beta\sigma_\gamma\right]\right)
		\label{eq:master_metric_preview}
	\end{equation}
	
	The unique diffeomorphism-invariant action on $\sigma$-space is:
	\begin{equation}
		\boxed{S[\sigma] = \int\!\dd^4x\,\sqrt{-g(\sigma)}\left[\frac{R[g(\sigma)]}{16\pi G} + \frac{\hbar^2}{2M}\,g^{\mu\nu}\nabla_\mu\sigma^\alpha\nabla_\nu\sigma_\alpha + \mathcal{L}_{\mathrm{eff}}(\psi,\,g(\sigma))\right]}
		\label{eq:qgd_action}
	\end{equation}
	
	Three terms:
	\begin{enumerate}
		\item \textbf{Einstein--Hilbert in $\sigma$-variables:} $R[g(\sigma)]/(16\pi G)$ is the unique second-order diffeomorphism-invariant scalar.  No independent gravitational postulate is introduced; the Einstein--Hilbert form is selected by requiring diffeomorphism invariance and at most second derivatives.
		\item \textbf{The $\sigma$-kinetic term:} $(\hbar^2/2M)(\nabla\sigma)^2$ is required by dimensional consistency ($\sigma$ is dimensionless; $\hbar^2/M$ is the unique available scale with dimensions of action$\times$mass$^{-1}$), stability, and its origin from coarse-graining the Dirac phase.  It is the gravitational analogue of an elastic stiffness modulus.
		\item \textbf{Matter coupling:} $\mathcal{L}_{\mathrm{eff}}$ (Eq.~\eqref{eq:Leff_final}) couples to the \emph{derived} metric $g(\sigma)$, preserving the equivalence principle.
	\end{enumerate}
	
	Varying with respect to $\sigma_\alpha$ (\textbf{not} $g_{\mu\nu}$) yields the fundamental field equation:
	\begin{equation}
		\boxed{\frac{\hbar^2}{M}\,\nabla^2\sigma^\alpha = \frac{1}{16\pi G}\!\left(G^{\mu\nu} - \frac{8\pi G}{c^4}T^{\mu\nu}\right)\frac{\partial g_{\mu\nu}}{\partial\sigma_\alpha}}
		\label{eq:sigma_field_eq}
	\end{equation}
	
	This equation has no direct counterpart in general relativity, where the metric is the fundamental dynamical variable and no wave equation for an underlying field is present.
	
	At equilibrium ($\nabla^2\sigma = 0$), the left-hand side vanishes, recovering:
	\begin{equation}
		\boxed{G_{\mu\nu} = \frac{8\pi G}{c^4}\,T_{\mu\nu}}
	\end{equation}
	
	The dichotomy is sharp:
	\begin{equation}
		\begin{cases}
			\nabla^2\sigma = 0 &\longrightarrow\quad\text{Classical GR (static $\sigma$ configuration)} \\[4pt]
			\nabla^2\sigma \neq 0 &\longrightarrow\quad\text{Quantum gravity ($\sigma$ fluctuations)}
		\end{cases}
	\end{equation}
	
	In this sense, GR may be interpreted as the equilibrium limit of the graviton field dynamics, in the same way that thermodynamics describes the equilibrium limit of statistical mechanics.

	%=============================================================
	\section{Conclusion}
	%=============================================================
	
	The present paper has developed a formulation of gravitational dynamics in which the wavefunction amplitude is connected to classical momentum via the current-conservation relation $|\psi|^2=C/|p|$, and in which the Newtonian potential emerges from matching the quantum phase-field energy to the gravitational potential of a point mass.  The construction is carried out entirely in flat Minkowski spacetime.
	
	The gravitational fine structure constant
	\begin{equation*}
		\alpha_G = e^{i\pi/4}\sqrt{\frac{c\hbar}{2GMm}}
	\end{equation*}
	is \emph{complex}---a feature, not a deficiency.  The phase $e^{i\pi/4}$ encodes the geometric phase of the gravitational interaction.  Physical observables use $|\alpha_G|^2$ (probability densities) or $\mathrm{Re}(\alpha_G)$ (real amplitudes).
	
	With this corrected coupling, the probability flux simplifies to $C=mc/\sqrt{2}$, showing the flux scale is set by the test particle's rest momentum alone.  All source-mass dependence enters through $\alpha_G$ and $\Delta$.
	
	The momentum structure satisfies the cubic $|p|+c^2|p|^3/\Delta^2=\mathcal{P}$, where $\mathcal{P}=\sqrt{2}\,GMm^2/\hbar$ is the fundamental gravitational momentum scale.  The graviton scalar $\sigma=pc/\Delta\sim\beta=v/c$ provides a dimensionless measure of quantum gravitational strength with cubic dispersion $\sigma+\sigma^3=\mathcal{P}c/\Delta$.
	
	All standard solutions (Schwarzschild, Kerr, Reissner-Nordstr\"om, cosmological fluids) emerge as special cases in flat Minkowski spacetime.
	
	\bigskip
	The action~\eqref{eq:qgd_action} governs the dynamics of $\sigma_\mu$; the stress-energy tensor sourcing it is derived from the same Lagrangian; and the complex wavefunction with the $1/r$ spherical-wave factor provides the underlying quantum-mechanical realisation of the gravitational field.

\end{document}
